\documentclass[12pt]{article}
\usepackage[margin=1in]{geometry}
\usepackage{amsmath,amssymb,amsthm}
\usepackage{hyperref}

\newtheorem{theorem}{Theorem}
\newtheorem{lemma}[theorem]{Lemma}
\newtheorem{proposition}[theorem]{Proposition}
\newtheorem{corollary}[theorem]{Corollary}
\theoremstyle{definition}
\newtheorem{definition}[theorem]{Definition}
\theoremstyle{remark}
\newtheorem{remark}[theorem]{Remark}

\newcommand{\R}{\mathbb{R}}
\newcommand{\C}{\mathbb{C}}
\newcommand{\Z}{\mathbb{Z}}
\newcommand{\E}{\mathbb{E}}
\newcommand{\Var}{\mathrm{Var}}
\newcommand{\Tr}{\mathrm{Tr}}
\newcommand{\ph}{\varphi}
\newcommand{\La}{\Lambda}
\newcommand{\im}{\mathrm{Im}}
\newcommand{\re}{\mathrm{Re}}

\title{Mass Gap for the O($N$) Model: Analysis of the HS Approach}
\author{M.~R.~Douglas}
\date{April 2026}

\begin{document}
\maketitle

\begin{abstract}
We analyze the Hubbard-Stratonovich approach to proving the mass
gap for the O($N$) model in two dimensions at large~$N$. The HS
identity for the Euclidean quartic requires imaginary coupling,
creating a sign problem. We compare the NLSM (where Kupiainen's
constraint trick gives a real mass $m_0^2$ in the operator, enabling
a direct FK bound for each $\sigma$-value) with the LSM (where the
real mass cancels after HS, and the mass gap must emerge from the
oscillatory $\sigma$-average). A comparison across $d=0,1,2$ shows
that the sign problem is specific to $d=2$: the logarithmic UV
divergence of the Green's function promotes the self-energy from
$O(1/N)$ to $O(1)$, causing dimensional transmutation. We propose
a proof strategy based on the \textbf{Lefschetz thimble}: a contour
deformation to the steepest descent manifold through the dominant
saddle, where the integrand becomes real, positive, and approximately
log-concave. At large~$N$, a single thimble dominates (exponential
suppression of competitors), eliminating the sign problem and
enabling Brascamp-Lieb concentration and Feynman-Kac bounds. We
compare four strategies and discuss the current status of the Lean
formalization.
\end{abstract}

\tableofcontents

\section{The HS identity and the sign problem}

\subsection{The Euclidean quartic requires imaginary coupling}

The O($N$) LSM action on $\La \subset \Z^2$:
\begin{equation}
  S[\ph] = \tfrac{1}{2}\sum_i\langle\ph^i,(-\Delta)\ph^i\rangle
    + N\sum_x\lambda\bigl(\tfrac{|\ph(x)|^2}{N}-\rho^2\bigr)^2.
\end{equation}
The Boltzmann weight is $e^{-S}$. The quartic appears as
$e^{-\lambda a^2}$ with $a = |\ph|^2/N - \rho^2$. The HS identity:
\begin{equation}
\label{eq:HS}
  e^{-\lambda a^2}
    = \frac{1}{\sqrt{4\pi\lambda}}\int_{-\infty}^{\infty}
      e^{-\sigma^2/(4\lambda)\;+\;i\sigma a}\,d\sigma.
\end{equation}
Here $\sigma \in \R$ (real integration variable),
$e^{-\sigma^2/(4\lambda)}$ ensures convergence, and the coupling
$i\sigma a$ is \textbf{imaginary}. This is a Fourier transform
(proved from Mathlib's \texttt{fourierIntegral\_gaussian}).

The imaginary coupling is forced by the \emph{minus} sign in
$e^{-\lambda a^2}$ (repulsive/confining quartic in Euclidean
signature). A real coupling would give $e^{+\lambda a^2}$
(attractive, wrong sign).

\subsection{The $\sigma$-integral on the real axis}

After applying~\eqref{eq:HS} at each site and integrating out $\ph$
(Gaussian with complex mass), the partition function becomes:
\begin{equation}
\label{eq:f}
  Z = c\!\int_{\R^{|\La|}}\!
    \det\bigl(-\Delta + 2i\sigma z\bigr)^{-N/2}\,
    e^{-\sum_x[\sigma(x)^2/(4\lambda)+i\sigma(x)\rho^2]}\,
    d^{|\La|}\sigma,
\end{equation}
where $z = N^{-1/2}$ and the $\det$ comes from the $\ph$-integral.

This integral \textbf{converges} on the real $\sigma$-axis:
\begin{itemize}
\item $e^{-\sigma^2/(4\lambda)}$ provides Gaussian decay in $\sigma$.
\item $|\det(-\Delta+2i\sigma z)^{-N/2}|
  \le \det(-\Delta)^{-N/2}$ (bounded, since $|1+i\theta| \ge 1$
  for real~$\theta$).
\end{itemize}
But the integrand is \textbf{complex} (the det and the $e^{i\sigma \rho^2}$
factor oscillate). This is the sign problem.

\section{The NLSM: Kupiainen's mass trick}

\subsection{Adding $m_0^2$ for free}

For the NLSM with constraint $|\ph|^2 = N\beta$: any function of
$|\ph|^2$ is constant on the constraint surface. In particular,
$\tfrac{1}{2}m_0^2|\ph|^2 = \tfrac{1}{2}m_0^2 N\beta = \text{const}$.

So we can \textbf{freely add} $\tfrac{1}{2}m_0^2|\ph|^2$ to the
action, giving the massive GFF:
\begin{equation}
  S_{\rm NLSM} = \tfrac{1}{2}\langle\ph,(-\Delta+m_0^2)\ph\rangle
    + \text{const},
  \qquad |\ph|^2 = N\beta.
\end{equation}
The parameter $m_0^2 > 0$ is arbitrary; we choose it via the gap
equation (Kupiainen eq.~9b):
$(-\Delta+m_0^2)^{-1}_{ii} = \beta$.

\subsection{The NLSM operator has real mass}

After HS (introducing the Lagrange multiplier $\alpha$ for the
constraint), the $\ph$-propagator at fixed $\alpha$ is:
\begin{equation}
  \bigl(-\Delta + m_0^2 - 2i\alpha z\bigr)^{-1}(x,0).
\end{equation}
The operator has \textbf{real part} $-\Delta + m_0^2 \ge m_0^2 > 0$.
By the triangle inequality (FK representation):
\begin{equation}
\label{eq:NLSMbound}
  \bigl|(-\Delta+m_0^2-2i\alpha z)^{-1}(x,0)\bigr|
    \le (-\Delta+m_0^2)^{-1}(x,0)
    \sim e^{-m_0|x|}.
\end{equation}
\textbf{Every} $\alpha$-value gives exponential decay with mass $m_0$.
The $\alpha$-average is a correction, not the leading effect.

\subsection{Kupiainen's proof structure}

\begin{enumerate}
\item Gap equation: $m_0 > 0$ (always in 2d).
\item Bound~\eqref{eq:NLSMbound}: each $\alpha$-value gives decay.
\item Cluster expansion + RP: control the $\alpha$-average
  (prefactor, not the mass).
\item Mass gap: $m = m_0(1 + O(1/N))$.
\end{enumerate}

The hard step is (3), but the mass $m_0$ comes for free from (2).

\section{The LSM: the mass cancels}

\subsection{Add-subtract fails}

For the LSM (no constraint), we try the same trick: add
$\tfrac{1}{2}m_0^2|\ph|^2$ to the action and subtract it from
the interaction:
\begin{equation}
  S = \bigl[\tfrac{1}{2}\langle\ph,(-\Delta+m_0^2)\ph\rangle\bigr]
    + \bigl[\lambda\bigl(\tfrac{|\ph|^2}{N}-\rho^2\bigr)^2
      - \tfrac{1}{2}m_0^2\,\tfrac{|\ph|^2}{N}\bigr].
\end{equation}
Apply HS to $\lambda(|\ph|^2/N-\rho^2)^2$ only. The subtracted mass
$-\tfrac{1}{2}m_0^2\,|\ph|^2/N$ remains outside the HS. Computing the
$\ph$-dependence:
\begin{align}
  &-\tfrac{1}{2}\langle\ph,(-\Delta+m_0^2)\ph\rangle
    + i\sigma|\ph|^2 + \tfrac{1}{2}m_0^2|\ph|^2 \nonumber\\
  &\quad= -\tfrac{1}{2}\langle\ph,(-\Delta)\ph\rangle + i\sigma|\ph|^2
    \nonumber\\
  &\quad= -\tfrac{1}{2}\langle\ph,(-\Delta - 2i\sigma)\ph\rangle.
\end{align}
The $m_0^2$ from the GFF and the $-m_0^2$ from the subtraction
\textbf{cancel exactly}. The LSM operator is:
\begin{equation}
\label{eq:LSMop}
  -\Delta - 2i\sigma \qquad\text{(\textbf{no real mass})}.
\end{equation}

\subsection{Consequences}

The operator~\eqref{eq:LSMop} has eigenvalues $k^2 - 2i\sigma$
with real part $k^2 \ge 0$ (from $-\Delta$) and imaginary part
$-2\sigma$. For $k = 0$: the eigenvalue is $-2i\sigma$ (purely
imaginary).

\begin{itemize}
\item \textbf{No single-$\sigma$ decay}: $|(-\Delta-2i\sigma)^{-1}(x,0)|$
  does NOT decay exponentially for each $\sigma$-value. At $k=0$:
  $|1/(-2i\sigma)| = 1/(2|\sigma|)$, independent of $|x|$.
\item \textbf{Mass from $\sigma$-average only}: The exponential decay
  of $\langle\ph(x)\ph(0)\rangle$ must come \emph{entirely} from the
  oscillatory cancellations in the $\sigma$-integral~\eqref{eq:f}.
\item \textbf{The FK bound fails}: The NLSM bound~\eqref{eq:NLSMbound}
  relies on the real mass $m_0^2$. Without it, $|G_\sigma| \le G_0$
  gives $G_0 = (-\Delta)^{-1}$ (massless propagator, no decay).
\end{itemize}

\subsection{Comparison}

\begin{center}
\begin{tabular}{|l|c|c|}
\hline
& NLSM & LSM \\
\hline
Mass in operator & $m_0^2$ (free, from constraint) & \textbf{none} (cancels) \\
Operator after HS & $-\Delta+m_0^2-2i\alpha z$ & $-\Delta-2i\sigma$ \\
Single-$\sigma$ decay & $e^{-m_0|x|}$ (yes) & \textbf{none} \\
Mass gap source & Real mass + corrections & \textbf{$\sigma$-average only} \\
FK bound & Works ($|G| \le G_{m_0}$) & \textbf{Fails} \\
\hline
\end{tabular}
\end{center}

\section{Comparison across dimensions}

The sign problem and mass gap structure depend strongly on dimension.

\subsection{$d = 0$ (matrix model)}

One site ($|\La|=1$). The $\sigma$-integral is one-dimensional:
$\int_{-\infty}^\infty e^{-\sigma^2/(4\lambda)+i\sigma a}\,d\sigma$.
The saddle-point Gaussian is exact up to $O(1/N)$ corrections.
Oscillations are $O(1/N^2)$ --- negligible. No sign problem.

The Green's function $C = 1/m_0^2$ is a number. The self-energy
$\Sigma = O(1/N)$. The mass is perturbative: $m = m_{\rm bare}(1+O(1/N))$.
No dimensional transmutation.

\subsection{$d = 1$ (quantum mechanics)}

$|\La| = L/a$ sites (discretized time). The Green's function
$C(t,t) = 1/(2m_0)$ is \textbf{finite} (no UV divergence).
The self-energy $\Sigma = O(1/N)$ (no log divergence).

Oscillations at the saddle width: $\sim L/(aN^2)$.
For fixed $L, a$ and large $N$: oscillations vanish.
The saddle-point expansion converges for $N > C\cdot L/a$.
The mass gap is perturbative.

\subsection{$d = 2$ (QFT)}

$|\La| = (L/a)^2$ sites. The Green's function
$C(x,x) = \frac{1}{4\pi}\log\frac{\La^2}{m_0^2}$ \textbf{diverges}
logarithmically with the UV cutoff.

The self-energy:
$\Sigma \sim \frac{1}{N}\cdot\log\frac{\La}{m_0} \sim \frac{1}{N}\cdot N = O(1)$.
The $1/N$ from BL is compensated by the log divergence ---
\textbf{dimensional transmutation}. The gap equation gives
$m_0 \sim \La\,e^{-c/g^2}$ (exponentially small).

Oscillations at the saddle width: $\sim |\La|/N^2 = L^2/(a^2 N^2)$.
For the continuum limit $a \to 0$: oscillations grow, sign problem
becomes severe.

\subsection{Summary}

\begin{center}
\begin{tabular}{|l|c|c|c|}
\hline
& $d=0$ & $d=1$ & $d=2$ \\
\hline
$C(x,x)$ & $1/m_0^2$ & $1/(2m_0)$ & $\frac{1}{4\pi}\log\frac{\La}{m_0}$ \\
UV divergence & none & none & \textbf{log} \\
Self-energy $\Sigma$ & $O(1/N)$ & $O(1/N)$ & $\mathbf{O(1)}$ \\
Gap equation & trivial & perturbative & \textbf{dim.\ transmutation} \\
Oscillation at width & $O(1/N^2)$ & $O(L/(aN^2))$ & $O(L^2/(a^2N^2))$ \\
Sign problem & none & mild & \textbf{severe} \\
\hline
\end{tabular}
\end{center}

The fundamental difference: in $d=2$, the log divergence of $C(x,x)$
makes the self-energy $O(1)$ instead of $O(1/N)$, producing dimensional
transmutation and making the sign problem unavoidable.

\section{The saddle-point structure}

Despite the difference in operator structure, both models share
the same saddle-point analysis.

\subsection{The gap equation}

On the real $\sigma$-axis, the saddle of the exponent in
\eqref{eq:f} is at $\sigma = 0$ (after absorbing the saddle
value into $m_0^2$ for the NLSM, or into the gap equation for the LSM).
The gap equation in both cases:
\begin{equation}
  (-\Delta+m_0^2)^{-1}_{xx} = \beta_{\rm eff}
\end{equation}
where $\beta_{\rm eff}$ encodes the coupling. In 2d:
$m_0 \sim \La\,e^{-c/g^2}$ (exponentially small).

\subsection{The saddle-point Gaussian}

Near $\sigma = 0$: the exponent is approximately
$-\tfrac{1}{2}(\sigma, H\sigma)$ where $H > 0$ (positive definite
Hessian, including both the bare $1/(2\lambda)$ and the $\Tr\log$
contribution). The saddle-point Gaussian IS log-concave on the
real axis.

\subsection{The oscillatory corrections}

Away from $\sigma = 0$: the exponent develops an imaginary part
$\mathrm{Im}\,f(\sigma) \ne 0$, causing oscillations. The sign problem
is the difficulty of controlling these oscillations.

Key scales (at the saddle):
\begin{itemize}
\item Width of the Gaussian: $\sigma \sim 1/\sqrt{N}$ per site.
\item Oscillation onset: $\mathrm{Im}\,f \sim O(\sigma^3/\sqrt{N})$
  (cubic, since $\mathrm{Im}\,f'' = 0$ at the saddle).
\item Total oscillation at the saddle width: $\sim |\La|/N^2$.
\end{itemize}
For $N \gg \sqrt{|\La|}$: oscillations are negligible within the
saddle-point width. The saddle-point approximation is accurate.

\section{What Kupiainen proved}

\subsection{NLSM (1980a): mass gap}

Kupiainen~\cite{Kupiainen80} proved the mass gap for the lattice
NLSM at large~$N$: exponential decay of correlations in the odd
sector, with mass $m(\beta,n) \to m_0(\beta)$ as $n \to \infty$.

The proof uses the constraint trick ($m_0^2$ added for free), the
dual representation with complex measure, reflection positivity,
chessboard estimates, and a random walk expansion. The threshold
is $n > a_1 m_0^{-a_2}$ with $a_2 = 6(d+2)+1 = 25$ in $d = 2$.

\subsection{LSM (1980b): pressure only}

Kupiainen~\cite{Kupiainen80b} studied the continuum $(\ph^2)_2^2$
model (the LSM) and proved a \textbf{weaker} result: the $1/n$
expansion for the \textbf{pressure} (vacuum energy) is asymptotic:
\begin{equation}
  p = \sum_{j=0}^{m-1} p_j n^{-j} + R_m,
  \qquad |R_m| \le (3m)!\,\lambda^4\alpha^m
    \Bigl(\frac{\lambda^\beta}{n}\Bigr)^m.
\end{equation}
The effective small parameter is $\lambda^\beta/n$ (with $\beta$
close to~10), not $e^{c\lambda}/n$.

\textbf{He did NOT prove the mass gap for the LSM.} The paper
works in the $\ph$-representation (integration by parts, no HS)
and in the dual representation (complex nonlocal measure, same
sign problem). He mentions that the mass gap would need ``a new
cluster expansion which explicitly cancels the diverging pressure
factors'' --- left as future work that was apparently never published.

\subsection{Implications}

\begin{itemize}
\item The \textbf{NLSM mass gap} at large $N$ is proved (Kupiainen 1980a).
\item The \textbf{LSM mass gap} at large $N$ is \textbf{open} ---
  only the pressure expansion is known (Kupiainen 1980b).
\item Proving the LSM mass gap would be a \textbf{genuinely new result},
  not a reformalization of existing work.
\item The LSM and NLSM share the same IR physics, so the mass gap
  ``should'' transfer. But making this rigorous is nontrivial.
\end{itemize}

\section{Strategy: Lefschetz thimble}\label{sec:thimble}

The right contour deformation is not a rotation to the imaginary
axis (which diverges), but the \textbf{Lefschetz thimble}: the
steepest descent manifold through the dominant saddle point.
This connects our problem to a well-developed mathematical
framework (Picard-Lefschetz theory) and a large recent literature
on sign problems in lattice QFT. See \texttt{sign-problem.tex}
for a general survey.

\subsection{Lefschetz thimbles: definition}

The exponent $f : \C^{|\La|} \to \C$ from~\eqref{eq:f} is
holomorphic (away from the singularities of the $\log\det$).
For each critical point $\sigma_*$ of~$f$, the
\textbf{Lefschetz thimble} $\mathcal{J}_{\sigma_*}$ is the
stable manifold of $\sigma_*$ under the anti-holomorphic
gradient flow:
\begin{equation}\label{eq:flow}
  \frac{d\sigma}{dt} = \overline{\nabla f(\sigma)}, \qquad
  \mathcal{J}_{\sigma_*} = \bigl\{\sigma(0) :
    \lim_{t\to+\infty}\sigma(t) = \sigma_*\bigr\}.
\end{equation}
Key properties:
\begin{enumerate}
\item $\mathcal{J}_{\sigma_*}$ is an $|\La|$-dimensional real
  submanifold of~$\C^{|\La|}$.
\item $\im\,f = \im\,f(\sigma_*)$ is \textbf{constant} on
  $\mathcal{J}_{\sigma_*}$ (the flow preserves the imaginary part).
\item $\re\,f$ is \textbf{maximized} at $\sigma_*$ on
  $\mathcal{J}_{\sigma_*}$ (steepest descent from the saddle).
\item The integrand $e^f$ restricted to $\mathcal{J}_{\sigma_*}$
  is proportional to a \textbf{positive} function (constant
  phase $e^{i\,\im f(\sigma_*)}$, positive amplitude
  $e^{\re f}$).
\end{enumerate}
Our ``steepest descent contour'' $\Gamma = \{\im f = 0,\;
\text{connected to } \sigma_* = 0\}$ is precisely the Lefschetz
thimble $\mathcal{J}_0$ for the dominant saddle at
$\sigma_* = 0$ (where $f(0)$ is real, so $\im f(\sigma_*) = 0$).

\subsection{Thimble decomposition}

By Picard-Lefschetz theory, the original integration cycle
decomposes as:
\begin{equation}\label{eq:decomp}
  [\R^{|\La|}] = \sum_{\sigma_*} n_{\sigma_*}\,
    [\mathcal{J}_{\sigma_*}],
\end{equation}
where $n_{\sigma_*} \in \Z$ are intersection numbers
(of $\R^{|\La|}$ with the dual ascending manifolds
$\mathcal{K}_{\sigma_*}$). Therefore:
\begin{equation}
  Z = \sum_{\sigma_*} n_{\sigma_*}\,e^{i\,\im f(\sigma_*)}\,
    Z_{\sigma_*}, \qquad
  Z_{\sigma_*} = \int_{\mathcal{J}_{\sigma_*}}
    e^{\re f(\sigma) - \re f(\sigma_*)}\,|d\sigma| > 0.
\end{equation}
Each $Z_{\sigma_*}$ is \textbf{real and positive}: the sign
problem is solved per-thimble. But different thimbles contribute
different phases $e^{i\,\im f(\sigma_*)}$, creating a residual
\emph{inter-thimble sign problem} if multiple thimbles contribute.

\subsection{Single-thimble dominance at large $N$}

The exponent scales as $f \sim N\cdot g(\sigma)$ for a function~$g$
independent of~$N$ (since $f$ contains the $N/2$ prefactor of
$\Tr\log$ and the $\sigma^2/(4\lambda)$ term scales with the
gap equation). Therefore:
\begin{itemize}
\item The Hessian scales as $f'' \sim N$: saddle-point width
  $\sim 1/\sqrt{N}$.
\item At competing saddle points $\sigma'_*$:
  $\re f(\sigma_*) - \re f(\sigma'_*) \sim N\cdot\Delta g > 0$
  (the action gap grows linearly in~$N$).
\item The dominant thimble contribution exceeds all others by
  $e^{cN}$ for some $c > 0$.
\end{itemize}

\textbf{Conclusion}: For $N$ sufficiently large, a single thimble
captures~$Z$ up to exponentially small corrections:
\begin{equation}
  Z = n_0\,Z_0\,\bigl(1 + O(e^{-cN})\bigr).
\end{equation}
This eliminates both the intra-thimble sign problem (the integrand
is positive on $\mathcal{J}_0$) and the inter-thimble sign problem
(other thimbles are exponentially suppressed).

\subsection{No topological obstruction for scalar fields}

Lawrence-Yamauchi~\cite{LY24} proved that for 2D Abelian lattice
gauge theory, the sign problem \emph{cannot} be removed by any
contour deformation. The obstruction is topological: $\pi_1(U(1))=\Z$
creates winding sectors with irreconcilable phases.

This obstruction does \textbf{not} apply to scalar field theories:
the field space $\R^n$ is contractible ($\pi_k(\R^n) = 0$ for
all~$k$), so there is no topological obstruction to deforming
the contour. For the O($N$) model, the question of whether
contour deformation solves the sign problem is \textbf{analytic}
(singularity avoidance, decay at infinity), not topological.

\section{The constant shift and thimble geometry}
\label{sec:constant_shift}

The thimble has a remarkably simple leading-order form: a constant
imaginary shift $v = v_*$. This section computes the thimble
explicitly, first in $d = 0$, then at general $|\La|$, and
establishes its key geometric property (Lagrangian submanifold).

\subsection{The $d = 0$ computation}

Consider one site ($|\La|=1$) with $N$ components of~$\ph$.
After HS and integrating out~$\ph$:
\begin{equation}\label{eq:f_d0}
  f(\sigma) = -\frac{\sigma^2}{4\lambda} + i\sigma \rho^2
    - \frac{N}{2}\log(i\sigma) + \text{const}.
\end{equation}
The saddle equation $f'(\sigma) = 0$ gives
$\sigma^2 - 2i\lambda \rho^2\sigma + \lambda N = 0$, with solutions:
\begin{equation}
  \sigma_* = i\bigl[\lambda \rho^2 \mp \sqrt{\lambda^2 \rho^4 + \lambda N}\,\bigr].
\end{equation}
Both saddle points are \textbf{purely imaginary}: $\sigma_* = iv_*$
with $v_* \in \R$.

Writing $\sigma = u + iv$ ($u$ = real part, $v$ = imaginary part),
$f(\sigma_*)$ is real (since $\sigma_*$ is purely imaginary), so
$\im f(\sigma_*) = 0$. The thimble equation $\im f = 0$ becomes:
\begin{equation}\label{eq:thimble_eq}
  u\Bigl[-\frac{v}{2\lambda} + \rho^2\Bigr]
    = \frac{N}{2}\arctan\!\Bigl(\frac{-u}{v}\Bigr).
\end{equation}
Expanding for small $u$: $\arctan(-u/v) \approx -u/v$, so
dividing by~$u$:
\begin{equation}
  -\frac{v}{2\lambda} + \rho^2 = -\frac{N}{2v}
  \quad\Longrightarrow\quad
  v^2 - 2\lambda \rho^2 v - \lambda N = 0,
\end{equation}
which gives $v = v_*$ (the saddle value). Near the saddle,
\textbf{the thimble is a horizontal line} $\sigma = u + iv_*$ ---
a constant imaginary shift.

\subsection{The operator on the shifted contour}

On the constant-shift contour $\sigma = u + iv_*$
(i.e., $v = v_*$ for all sites), with $v_*$ chosen via the
gap equation so that $-2v_* z = m_0^2$, the $\ph$-operator becomes:
\begin{equation}\label{eq:shifted_op}
  -\Delta + 2i\sigma z = -\Delta + 2i(u + iv_*)z
    = -\Delta + m_0^2 + 2iuz.
\end{equation}
(We write $\alpha \equiv v_*$ when convenient.)
The constant shift introduces the \textbf{real mass}~$m_0^2$ into
the operator. For each real field $u(x)$, the operator
$-\Delta + m_0^2 + 2iuz$ has positive real part $-\Delta + m_0^2$,
and the FK bound gives:
\begin{equation}
  \bigl|(-\Delta + m_0^2 + 2iuz)^{-1}(x,0)\bigr|
    \le (-\Delta + m_0^2)^{-1}(x,0)
    \sim e^{-m_0|x|}.
\end{equation}
\textbf{This is exactly Kupiainen's NLSM bound}, achieved by
contour deformation rather than by the constraint.

\subsection{The thimble unifies LSM and NLSM}

The constant imaginary shift has a conceptual interpretation:
\begin{itemize}
\item For the \textbf{NLSM}: the constraint $|\ph|^2 = N\beta$
  allows adding $m_0^2$ for free. No contour deformation needed.
\item For the \textbf{LSM}: the add-subtract trick cancels $m_0^2$
  on the real $\sigma$-axis ($v = 0$). But the \emph{thimble} passes
  through $\sigma_* = iv_*$, and on the thimble, $m_0^2$ reappears
  as the real part $-2v_* z$.
\item The contour deformation $v = 0 \to v = v_*$
  \textbf{does for the LSM what the constraint does for the NLSM}:
  it introduces a real mass into the $\ph$-operator.
\end{itemize}

\subsection{The $N = \infty$ thimble is exactly flat}

At the Gaussian (quadratic) level, the Hessian at the saddle is
$H = -f''(\sigma_*)$, which is \textbf{real and positive definite}
(since the operator at the saddle is $-\Delta + m_0^2$, real).
Writing $\sigma = u + iv$ and expanding around $\sigma_* = iv_*$
(so $\delta u = u$, $\delta v = v - v_*$):
\begin{equation}
  \im f = -(\delta u,\, H\,\delta v) \quad\text{(at quadratic order)}.
\end{equation}
Setting $\im f = 0$ with $H > 0$: this forces $\delta v = 0$,
i.e., $v = v_*$ everywhere.
The $N = \infty$ thimble is the hyperplane
$\sigma(x) = u(x) + iv_*$, \textbf{exactly flat}.

Since the Gaussian approximation is exact at $N = \infty$, the
thimble is exactly the constant shift at leading order. Corrections
come from the cubic and higher terms in~$f$, suppressed by
$1/\sqrt{N}$.

\subsection{The leading $1/N$ correction: explicit formula}
\label{sec:1N_correction}

The cubic term of $f$ at the saddle determines the first correction.
Expanding $\Tr\log(A_0 + 2i\delta\sigma z)$ in powers of
$\delta\sigma$:
\begin{equation}
  f'''_{xyz} = \frac{8i}{\sqrt{N}}\,G_0(x-y)\,G_0(y-z)\,G_0(z-x),
\end{equation}
where $G_0(k) = (k^2 + m_0^2)^{-1}$ is the massive propagator at
the saddle. The cubic coefficient is \textbf{purely imaginary}
(the factor $i$ comes from $(2iz)^3 = -8i/N^{3/2}$, and $G_0$ is
real at the saddle). Writing $f''' = iR$ with
$R_{xyz} = (8/\sqrt{N})\,G_0(x{-}y)\,G_0(y{-}z)\,G_0(z{-}x)$
(the \textbf{triangle diagram}).

The thimble $v = v(u)$ is parameterized as $v_i = \tfrac{1}{2}
\sum_{jk} S_{ijk}\,u_j u_k$ (quadratic in~$u$, $S$ symmetric in
$j,k$). The stable manifold invariance condition
$D\Phi(u)\cdot\dot{u} = \dot{v}|_{v=\Phi(u)}$ (with
$\dot{u} = -Hu$ at leading order, $\dot{v} = Hv - \tfrac{1}{2}R[u,u,\cdot]$)
gives a linear equation for~$S$:
\begin{equation}\label{eq:stable_manifold}
  S(Hw_1, w_2) + S(w_1, Hw_2) + H\cdot S(w_1, w_2) = R(w_1, w_2)
  \qquad \text{for all } w_1, w_2 \in \R^{|\La|},
\end{equation}
where the first two terms have $H$ acting on an \emph{input} of~$S$,
and the third has $H$ acting on the \emph{output}.

In Fourier space --- with $p, q$ the lattice momenta of the
two $u$-inputs and $p{+}q$ the momentum of the $v$-output ---
equation~\eqref{eq:stable_manifold} becomes algebraic:
$\hat{S}(p,q)\,[\hat{H}(p) + \hat{H}(q) + \hat{H}(p{+}q)]
= \hat{R}(p,q)$, giving:
\begin{equation}\label{eq:S_fourier}
  \boxed{\hat{S}(p,q) = \frac{\hat{R}(p,q)}
    {\hat{H}(p) + \hat{H}(q) + \hat{H}(p+q)}}
\end{equation}
The denominator is the sum of three $\sigma$ inverse
propagators at the three legs of the triangle vertex.
The total symmetry in $p, q, -(p{+}q)$ confirms the Lagrangian
property (Section~\ref{sec:lagrangian}). The ingredients:
\begin{align}
  \hat{H}(k) &= \frac{1}{2\lambda} + 2B(k),
    &B(k) &= \frac{1}{|\La|}\sum_l G_0(l)\,G_0(k{-}l)
    \quad\text{(bubble)}, \\
  \hat{R}(p,q) &= \frac{8}{\sqrt{N}}\,T(p,q),
    &T(p,q) &= \frac{1}{|\La|}\sum_l G_0(l)\,G_0(p{-}l)\,G_0(q{+}l)
    \quad\text{(triangle)}.
\end{align}
The explicit thimble correction in Fourier space:
\begin{equation}\label{eq:thimble_1N}
  \hat{v}(k) = \frac{4}{\sqrt{N}\,|\La|}
    \sum_{p+q=k}
    \frac{T(p,q)}{\hat{H}(p)+\hat{H}(q)+\hat{H}(p{+}q)}\;
    \hat{u}(p)\,\hat{u}(q).
\end{equation}

\subsection{Properties of the correction}

\begin{enumerate}
\item \textbf{Standard Feynman diagram}:
  The kernel $T/(\hat{H}+\hat{H}+\hat{H})$ is the triangle
  vertex dressed by three $\sigma$-propagators $\hat{H}^{-1}$.
  This is a standard one-loop diagram in the $1/N$ expansion ---
  the thimble shape \emph{is} the diagrammatic expansion.

\item \textbf{No singularities}: The denominator
  $\hat{H}(p) + \hat{H}(q) + \hat{H}(p{+}q)$ is strictly positive
  (each $\hat{H}(k) > 0$), so $\hat{S}$ is smooth.

\item \textbf{Locality}: In position space, the kernel
  $K(x{-}y, x{-}z) \sim e^{-\min(m_0,M_\sigma)(|x-y|+|x-z|)}$
  decays exponentially. The thimble correction at site~$x$
  depends mainly on $u(y)$ for $|y-x| \lesssim 1/m_0$.

\item \textbf{Ultra-local limit}: For zero momentum ($p = q = k = 0$):
  $\hat{S}(0,0) = \hat{R}(0,0)/(3\hat{H}(0))$, recovering the
  $d = 0$ result $v = u^2/(6\hat{H}(0)\sqrt{N})$.

\item \textbf{Size at the saddle width}: The fluctuation
  $u \sim 1/\sqrt{N}$, so:
  \begin{center}
  \begin{tabular}{|l|c|}
  \hline
  Quantity & Scale \\
  \hline
  $u(x)$ (real fluctuation) & $1/\sqrt{N}$ \\
  $v(x)$ (thimble correction) & $1/N^{3/2}$ \\
  $v/u$ ratio & $1/N$ \\
  Residual $\im f$ on constant shift & $u^3/\sqrt{N} \sim 1/N^2$ per site \\
  \hline
  \end{tabular}
  \end{center}
  The thimble is very nearly flat within the Gaussian peak.
\end{enumerate}

\subsection{Multi-site: constant shift plus local corrections}

For the multi-site case ($|\La| > 1$), the full thimble
$\sigma = u + iv$ has:
\begin{equation}
  v(x) = v_* +
    \frac{4}{\sqrt{N}}\sum_{y,z} K(x{-}y,\,x{-}z)\,u(y)\,u(z)
    + O(1/N),
\end{equation}
where $K$ is the position-space kernel of $T/(\hat{H}+\hat{H}+\hat{H})$.
The leading-order form is the \textbf{constant shift}
$v = v_*$ (i.e., $\sigma = u + iv_*$); the correction is local,
quadratic in~$u$, and suppressed by $1/\sqrt{N}$.

\subsection{The thimble is a Lagrangian submanifold}
\label{sec:lagrangian}

The thimble $\mathcal{J}_0$ is a \textbf{Lagrangian submanifold}
of $(\C^{|\La|},\omega)$, where
$\omega = \sum_j du_j \wedge dv_j$ is the standard symplectic form
on $\C^{|\La|} \cong \R^{2|\La|}$.

\subsubsection{Proof}

The anti-holomorphic gradient flow~\eqref{eq:flow}, written in
real coordinates $\sigma = u + iv$, is:
\begin{equation}
  \frac{du}{dt} = \frac{\partial(\re f)}{\partial u}
    = \frac{\partial(\im f)}{\partial v}, \qquad
  \frac{dv}{dt} = \frac{\partial(\re f)}{\partial v}
    = -\frac{\partial(\im f)}{\partial u},
\end{equation}
where the second equalities are Cauchy-Riemann. This is the
\textbf{Hamiltonian flow} of $\im f$ with respect to~$\omega$.
Hamiltonian flows preserve the symplectic form.

Let $X, Y$ be tangent vectors to $\mathcal{J}_0$ at a point~$\sigma$.
Under the flow, they remain tangent to $\mathcal{J}_0$
(flow-invariance of the stable manifold). As $t \to +\infty$,
both flow into the stable subspace at~$\sigma_*$, which is
$\{v = 0\}$ (the real hyperplane). Since
$\omega|_{\{v=0\}} = 0$:
\begin{equation}
  \omega(X,Y) = \omega(X(t),Y(t)) \;\xrightarrow{t\to\infty}\; 0.
\end{equation}
But $\omega(X(t),Y(t))$ is constant (Hamiltonian flow preserves
$\omega$), so $\omega(X,Y) = 0$. \qed

\subsubsection{Consequences}

\begin{enumerate}
\item \textbf{Generating functional}: Since $\mathcal{J}_0$ is
  Lagrangian, the thimble correction satisfies
  $v = \nabla_u \Phi$ for a single scalar functional
  $\Phi : \R^{|\La|} \to \R$. The entire thimble is determined
  by~$\Phi$.

\item \textbf{Jacobian phase}: The induced volume form on a
  Lagrangian graph $\{(u, \nabla\Phi(u))\}$ has Jacobian
  $\det(I + i\nabla^2\Phi)$. Since $\nabla^2\Phi$ is symmetric
  with real eigenvalues~$\lambda_j$:
  \begin{equation}
    |\det J| = \prod_j\sqrt{1+\lambda_j^2}, \qquad
    \arg(\det J) = \sum_j \arctan(\lambda_j)
      = \Tr\arctan(\nabla^2\Phi).
  \end{equation}
  The Jacobian phase is \textbf{not zero} in general: a real
  symmetric matrix does not have $\pm$-paired eigenvalues.
  For $\nabla^2\Phi \sim u/\sqrt{N}$
  (from the leading thimble correction):
  $\arg(\det J) \approx \Tr(\nabla^2\Phi) \sim |\La|/N$.
  See Section~\ref{sec:quantum_HJ} for how to handle this.

\item \textbf{Integrability of the $1/N$ correction}: The total
  symmetry of $\hat{S}(p,q)$ in all three momenta $p, q, -(p{+}q)$
  (from~\eqref{eq:S_fourier}) is precisely the condition
  $\partial v_i/\partial u_j = \partial v_j/\partial u_i$, i.e.,
  the existence of~$\Phi$.
\end{enumerate}

\section{The Hamilton-Jacobi equation}
\label{sec:HJ}

The Lagrangian property reduces the thimble to a single scalar
functional $\Phi(u)$. This section derives the exact equation
for~$\Phi$, its perturbative expansion, and the convergence theory.

\subsection{The equation}

The condition $\im f = 0$ on the thimble $v = \nabla\Phi$ gives a
\textbf{single scalar equation} for a single scalar functional:
\begin{equation}\label{eq:HJ}
  \boxed{\im\,f\bigl(u + i\nabla_u\Phi(u)\bigr) = 0
    \qquad\text{for all } u \in \R^{|\La|}.}
\end{equation}
This is a Hamilton-Jacobi equation: $u$ is the ``position,''
$\nabla\Phi$ is the ``momentum,'' $\im f$ is the ``Hamiltonian,''
and $\Phi$ is the ``action'' (generating function of the
Lagrangian submanifold).

\subsection{Concrete form: the arctan equation}

Writing $\psi(x) = \partial\Phi/\partial u(x)$ for the thimble
correction field, the imaginary part is $v(x) = v_* + \psi(x)$
and $\sigma = u + iv$. The $\ph$-operator on the thimble is:
\begin{equation}
  -\Delta + 2i\sigma z
    = \underbrace{(-\Delta + m_0^2 - \tfrac{2\psi}{\sqrt{N}})}_{M(\psi)}
    \;+\; i\underbrace{\tfrac{2u}{\sqrt{N}}}_{V(u)},
\end{equation}
where $M(\psi)$ is a real Schr\"odinger operator (with potential
$m_0^2 - 2\psi/\sqrt{N}$) and $V(u)$ is a real multiplication operator.
Assuming $M > 0$ (the thimble stays above the spectral gap):
\begin{equation}
  \im\,\Tr\log(M + iV) = \Tr\arctan\bigl(M^{-1/2}\,V\,M^{-1/2}\bigr),
\end{equation}
where $\arctan$ acts on the operator $D = M^{-1/2}VM^{-1/2}$
(self-adjoint, real eigenvalues $d_k$) by
$\arctan(D) = \sum_k \arctan(d_k)\,|e_k\rangle\langle e_k|$.

The exact HJ equation~\eqref{eq:HJ} becomes:
\begin{equation}\label{eq:HJ_arctan}
  \boxed{\frac{1}{2\lambda}\!\sum_x u(x)\bigl(v_*{+}\psi(x)\bigr)
    + \frac{N}{2}\,\Tr\arctan\!\Bigl(
      M(\psi)^{-1/2}\,\frac{2u}{\sqrt{N}}\,M(\psi)^{-1/2}\Bigr)
    + \rho^2\!\sum_x u(x) = 0}
\end{equation}
for all $u \in \R^{|\La|}$, with $\psi = \nabla\Phi(u)$ and
$M(\psi) = -\Delta + m_0^2 - 2\psi/\sqrt{N}$.

\begin{remark}[Structure of the equation]
Equation~\eqref{eq:HJ_arctan} is \textbf{exact at finite~$N$} ---
no expansion in $1/N$. The three terms have clear roles:
\begin{itemize}
\item $\frac{1}{2\lambda}(u,v_*{+}\psi)$: the bare Gaussian
  contribution to $\im f$.
\item $\frac{N}{2}\Tr\arctan(\cdots)$: the determinantal
  contribution from integrating out~$\ph$. The $\arctan$ is the
  \textbf{phase} of $\det(M+iV)^{1/2}$ --- the argument of
  the complex determinant.
\item $\rho^2\sum u$: the linear coupling to the HS source.
\end{itemize}
At $\psi = 0$ (constant shift) and $u = 0$: the equation reduces to
the \textbf{gap equation} for~$m_0$ (the $\arctan$ linearizes to
$M_0^{-1}V$ and the trace gives $G_0(0)\sum u$).
\end{remark}

\subsection{Perturbative expansion}

Expanding $\Phi = \Phi_3 + \Phi_4 + \cdots$ in powers of~$u$
(equivalently, $1/\sqrt{N}$):
\begin{align}
  \Phi_3(u) &= \frac{1}{6}\sum_{ijk} S_{ijk}\,u_i u_j u_k
    &\text{(triangle diagram, Section~\ref{sec:1N_correction})}, \\
  \Phi_4(u) &= \frac{1}{24}\sum_{ijkl} S^{(4)}_{ijkl}\,
    u_i u_j u_k u_l
    &\text{(box diagram)}, \\
  &\;\;\vdots \nonumber
\end{align}
There is no $\Phi_1$ or $\Phi_2$: $\Phi_1 = 0$ by the gap equation
(the linear term in~\eqref{eq:HJ_arctan} vanishes at $u = 0$),
and $\Phi_2 = 0$ because the Hessian $f''(\sigma_*)$ is real (the
quadratic term in $\im f$ vanishes identically).
Each $\Phi_n$ for $n \ge 3$ is determined by matching order $u^n$
in~\eqref{eq:HJ_arctan}, giving $S^{(n)}$ in terms of $n$-loop
Feynman diagrams and lower-order data.

\subsection{Convergence of the $1/N$ expansion}

\begin{proposition}[Analyticity of the thimble]
For $|\La| < \infty$ and $m_0 > 0$, the generating functional
$\Phi(u)$ is real-analytic in $u \in \R^{|\La|}$.
The Taylor series $\Phi = \sum_{n \ge 3} \Phi_n$ converges
in the \textbf{sup-norm} $\|u\|_\infty < r$ with radius:
\begin{equation}\label{eq:radius}
  r \;\ge\; c\,m_0\,N^{1/4},
\end{equation}
where $c > 0$ depends on $\lambda$ but \textbf{not} on $|\La|$
or~$N$.
\end{proposition}

\begin{proof}[Proof sketch]
The thimble $\mathcal{J}_0$ is the stable manifold of~$\sigma_*$
under the anti-holomorphic gradient flow~\eqref{eq:flow}. Since
$f$ is holomorphic and the saddle is non-degenerate ($H > 0$),
the \textbf{analytic stable manifold theorem} guarantees that
$\Phi$ is real-analytic in~$u$.

The radius is limited by the nearest singularity of~$\Phi$,
which occurs when $M(\psi) = -\Delta + m_0^2 - 2\psi/\sqrt{N}$ hits
zero eigenvalue. This requires $|\psi(x)| \ge m_0^2\sqrt{N}/2$
\emph{at some site}~$x$ --- a \textbf{pointwise} (sup-norm)
condition, independent of~$|\La|$. Since $p \sim u^2$, the
singularity is at $\|u\|_\infty \sim m_0 N^{1/4}$.
\end{proof}

\subsection{The Jacobian phase and the order of limits}
\label{sec:order_of_limits}

\subsubsection{The residual phase problem}

On any contour $\sigma = u + i\nabla\Phi(u)$, the integral
acquires a Jacobian:
\begin{equation}\label{eq:jacobian_phase}
  \int_{\mathcal{J}_0} e^f\,d\sigma
    = \int_{\R^{|\La|}} e^{f(u+i\nabla\Phi)}\,
      \det(I + i\nabla^2\Phi)\,du.
\end{equation}
The total phase of the integrand is:
\begin{equation}
  \text{total phase} = \underbrace{\im f(u+i\nabla\Phi)}_{\text{action phase}}
    \;+\; \underbrace{\Tr\arctan(\nabla^2\Phi)}_{\text{Jacobian phase}}.
\end{equation}
The classical HJ equation~\eqref{eq:HJ_arctan} sets the action
phase to zero. But the Jacobian phase
$\Tr\arctan(\nabla^2\Phi) \ne 0$: a real symmetric matrix does
\emph{not} have $\pm$-paired eigenvalues, so $\arg\det J \ne 0$.

For the leading correction $\nabla^2\Phi \sim u/\sqrt{N}$:
$\Tr\arctan(\nabla^2\Phi) \approx \Tr(\nabla^2\Phi)
\sim |\La|/N$, with variance $\sim |\La|/N^2$.

\begin{center}
\begin{tabular}{|l|c|c|}
\hline
Contour & Action phase & Jacobian phase \\
\hline
Constant shift ($v = v_*$) & $O(|\La|/N^2)$ & 0 (flat) \\
Classical thimble ($\im f = 0$) & 0 & $O(|\La|/N^2)$ \\
\hline
\end{tabular}
\end{center}
Both have \textbf{the same total phase} $O(|\La|/N^2)$. The
classical thimble moves the phase from the action to the Jacobian
but does not eliminate it.

\subsubsection{Consequence for the mass gap}

The triangle inequality gives:
$|Z^{-1}\!\int G_\sigma e^f\,d\sigma|
  \le G_{m_0}(x,0) \cdot (\text{average sign})^{-1}$,
where the average sign $\sim e^{-c|\La|/N^2}$. For
$|\La| \to \infty$ at fixed~$N$, this bound diverges. The
contour shift (or classical thimble) alone requires
$N \gg \sqrt{|\La|}$.

\subsubsection{The quantum Hamilton-Jacobi equation}
\label{sec:quantum_HJ}

The Jacobian phase $\Tr\arctan(\nabla^2\Phi)$ is the
\textbf{residual sign problem} identified in the lattice
QFT literature~\cite{CRS12,TNK15,Scorzato15,Fukuma20}:
even on the steepest-descent thimble, the measure $\det J$
introduces a complex phase that must be handled. Previous
approaches treat this phase as a computational obstacle
(reweighting, worldvolume methods). We instead treat it as a
\textbf{solvable equation}: by adjusting $\Phi$ to cancel both
the action phase and the Jacobian phase simultaneously.

To eliminate \emph{both} phases, we must solve the
\textbf{quantum HJ equation}:
\begin{equation}\label{eq:quantum_HJ}
  \boxed{\im\,f\bigl(u + i\nabla\Phi\bigr)
    + \Tr\arctan(\nabla^2\Phi) = 0
    \qquad\text{for all } u \in \R^{|\La|}.}
\end{equation}
This incorporates the Jacobian measure into the phase condition.
The integrand in~\eqref{eq:jacobian_phase} then has
\textbf{exactly zero total phase}: $e^{f}\det J$ is real and
positive on the quantum thimble.

Equation~\eqref{eq:quantum_HJ} is a \textbf{second-order} PDE
in~$\Phi$ (due to the $\nabla^2\Phi$ in the $\Tr\arctan$), in
contrast to the first-order classical HJ
equation~\eqref{eq:HJ_arctan}. Existence of solutions is more
subtle: the classical thimble provides a natural starting point,
with the quantum correction $\delta\Phi = O(1/N)$ absorbing
the Jacobian phase.

The quantum HJ equation has the same perturbative structure:
$\Phi = \Phi_3 + \Phi_4 + \cdots$ with each order modified by
$O(1/N)$ relative to the classical HJ solution.

\subsubsection{Bounds-only approach to the quantum thimble}
\label{sec:quantum_bounds}

Just as for the classical HJ (Section~\ref{sec:HJ_solution}),
we do not need to \emph{solve} the quantum HJ equation ---
we need \textbf{existence} of a solution~$\Phi$ with
\textbf{a priori bounds}. On the quantum thimble, the full
integrand $e^f \det J$ is real and positive (both phases cancel
by construction). Define the \textbf{effective potential}:
\begin{equation}
  V_{\rm eff}(u) = -\re\,f\bigl(u + i\nabla\Phi(u)\bigr)
    - \log|\det(I + i\nabla^2\Phi(u))|.
\end{equation}
The BL measure is $e^{-V_{\rm eff}(u)}\,du$ --- real, positive,
no residual phase. The three a priori bounds needed are:

\begin{enumerate}
\item \textbf{Proximity}: $|\psi(x)| \le C_1\|u\|_\infty^2/\sqrt{N}$
  (same as classical --- the quantum correction is $O(1/N)$ smaller).
\item \textbf{Spectral gap}: $M(\psi) \ge m_0^2/2 > 0$
  (same bound, same proof).
\item \textbf{Hessian of $V_{\rm eff}$}:
  $\nabla^2 V_{\rm eff} \ge \kappa N$. At the constant shift
  ($\Phi=0$), the Hessian is (in Fourier space):
  \begin{equation}
    \hat{H}(k) = \frac{1}{2\lambda} + 2B(k), \quad
    B(k) = \frac{1}{|\La|}\sum_l G_0(l)\,G_0(k{-}l) \ge 0,
  \end{equation}
  which gives $H \ge 1/(2\lambda) > 0$. The $\Tr\log$ contributes
  $+2B(k) \sim N \cdot O(1/N) = O(1)$ per site (the bubble diagram
  is $O(1)$ in 2d), while the bare term scales as
  $1/(2\lambda) \sim 1/(c/N) = N/c$ for the LSM coupling $\lambda
  \sim c/N$. So $H \ge \kappa N$ with $\kappa = 1/(2c)$.

  On the full thimble ($\Phi \ne 0$): $V_{\rm eff}$ includes
  $-\log|\det(I+i\nabla^2\Phi)|$. Using the proved Hessian formula
  (the proved \texttt{hessian\_log\_det} theorem):
  \begin{equation}
    \frac{\partial^2}{\partial u_x\,\partial u_y}
      \log|\det(I+iJ)|_{u=0}
    = \Tr(S_x\,S_y)
    = \sum_{z,w} S_{zwx}\,S_{wzy}
  \end{equation}
  where $S_x$ denotes the slice of the 3-tensor~$S$
  from~\eqref{eq:S_fourier} at index~$x$, and the formula
  follows from $d^2(\log\det A)(H,K) = -\Tr(A^{-1}HA^{-1}K)$
  evaluated at $A = I$, $H = iS_x$, $K = iS_y$.

  In Fourier: $|\hat{S}(p,q)| \le 16c\,|T(p,q)|/(3N^{3/2})$
  (from $|\hat{R}| = (8/\sqrt{N})|T|$ and
  $\hat{H} \ge N/(2c)$). So $\Tr(S_x S_y)$ is $O(1/N^3)$
  per momentum mode --- negligible compared to the $O(N)$ bare
  Hessian. For $N$ sufficiently large, $H \ge \kappa N/2 > 0$.
\end{enumerate}

The quantum $\Phi$ differs from the classical $\Phi$ by
$O(1/N)$, so the bounds are the same to leading order. The
key advantage: \textbf{zero total phase} means the average
sign is exactly~1, and the mass gap bound is
\textbf{uniform in~$|\La|$} without any $N \gg \sqrt{|\La|}$
condition.

\begin{proposition}[Existence of the quantum thimble]
\label{prop:quantum_thimble}
Define $F(\Phi)(u) = \im f(u + i\nabla\Phi) + \Tr\arctan(\nabla^2\Phi)$.
For $N > N_0(m_0,\lambda)$ on a finite lattice~$\La$:
\begin{enumerate}
\item $F(0)(u) = O(\|u\|_\infty^3/\sqrt{N})$ per site
  (the residual phase at the constant shift).
\item The linearization $DF(0)$ is invertible: the leading term
  is $\partial(\im f)/\partial\psi \cdot \nabla(\cdot) +
  \Delta_u(\cdot)$, which is controlled by the Hessian $H > 0$.
\item By the \textbf{implicit function theorem}: there exists
  $\Phi$ with $F(\Phi) = 0$ and
  $\|\nabla\Phi\|_\infty \le C\,\|u\|_\infty^2/\sqrt{N}$
  (the proximity bound).
\end{enumerate}
\end{proposition}

\begin{proof}[Proof sketch]
$F$ is smooth (from the analyticity of $f$ and $\arctan$).

\textbf{Residual phase at $\Phi=0$}: $F(0)(u) = \im f(u+iv_*)$
since $\Tr\arctan(0) = 0$. The expansion of $\im f$ in powers
of~$u$: the linear term is $-\sum_x (u_x/2\lambda)\,v_*$ which
is real (contributes to $\re f$, not $\im f$). The quadratic
term $\im(u,Hu) = 0$ since $H$ is real. The leading imaginary
term is cubic: $\im f \sim \frac{i}{6}R[u,u,u]$ where
$R_{xyz} = (8/\sqrt{N})T(x,y,z)$ (the triangle diagram).
Per site: $|F(0)(u)| \lesssim \|u\|_\infty^3/\sqrt{N}$.

\textbf{Linearization}: $DF(0) \cdot \delta\Phi$ maps $\delta\Phi$
to $\im[\partial f/\partial v \cdot \nabla(\delta\Phi)]
+ \Tr(\nabla^2(\delta\Phi)/(1+0))$. In Fourier space:
$\widehat{DF(0)}(k) = \hat{H}(k)$ (the same Hessian as
in~\eqref{eq:S_fourier}). Since $\hat{H}(k) \ge 1/(2\lambda)
\sim N/c > 0$, $DF(0)$ is invertible with $\|DF(0)^{-1}\| \le c/N$.

\textbf{IFT}: Since $\|F(0)(u)\| = O(\|u\|^3/\sqrt{N})$ and
$\|DF(0)^{-1}\| \le c/N$, the quantitative IFT gives a solution
$\Phi$ with $\|\nabla\Phi\| \le (c/N) \cdot O(\|u\|^3/\sqrt{N})
= O(\|u\|^3\,c/(N\sqrt{N}))$. In the BL-relevant region
$\|u\| \lesssim 1/\sqrt{\kappa N}$: $|\nabla\Phi| \lesssim
c/(\kappa^{3/2} N^{5/2}) \ll m_0^2\sqrt{N}/2$ (singularity distance).
\end{proof}

\begin{remark}[Contour deformation validity]
\label{rem:contour_validity}
For Cauchy's theorem to justify $\int_{\R^{|\La|}} = \int_{\rm thimble}$:
\begin{enumerate}
\item \textbf{Holomorphy}: The integrand $e^f$ is holomorphic in each
  $\sigma_x$ away from the zeros of $\det(-\Delta+2i\sigma z)$.
  Singularity avoidance: in the BL-relevant region
  ($|u| \lesssim 1/\sqrt{\kappa N}$), the thimble displacement
  $|\psi| \lesssim c/(\kappa^{3/2}N^{5/2})$ is exponentially
  smaller than the singularity distance~$m_0^2\sqrt{N}/2$.
\item \textbf{Decay on connecting walls}: At $|u_x| \to \infty$
  with other $u_y$ fixed, the Gaussian factor $e^{-u_x^2/(4\lambda)}$
  dominates the $\Tr\log$ growth (which is polynomial in~$u_x$:
  $|\Tr\log(-\Delta+m_0^2+2iuz)| \le C_1 + C_2\log(1+|u|)$ per site).
  So the integrand decays as $e^{-u_x^2/(4\lambda)+O(\log|u|)}
  \to 0$ on the side walls. This justifies the single-variable
  contour shift (\texttt{vertical\_contour\_shift}, proved);
  the multi-variable version follows by Fubini applied $|\La|$ times.
\item \textbf{Multi-variable Cauchy}: The proved
  \texttt{vertical\_contour\_shift} shifts one variable at a time.
  By Fubini (iterated integrals over a product), the $|\La|$-fold
  composition shifts all variables simultaneously. Each step uses
  the holomorphy and decay from~(1)--(2).
\end{enumerate}
\end{remark}

\begin{remark}[IFT domain size and global control]
\label{rem:ift_domain}
The quantitative IFT gives a solution $\Phi$ in a ball of radius
$r \sim \|DF(0)^{-1}\|^{-1}/\|D^2F\|$ in $u$-space.
Since $\|DF(0)^{-1}\| \le c/N$ (from $H \ge N/(2c)$) and
$\|D^2F\| = O(1)$ (the second derivative of the phase equation
is bounded independently of~$N$), the IFT ball has radius
$r \sim N/c = O(N)$ in $\|u\|_\infty$.

The BL measure $e^{-V_{\rm eff}}$ is concentrated in the region
$\|u\|_\infty \lesssim 1/\sqrt{\kappa N} = O(N^{-1/2})$
(from the Hessian $\nabla^2 V_{\rm eff} \ge \kappa N$).

The comparison $R_{\rm IFT} \sim O(N) \gg R_{\rm BL} \sim O(N^{-1/2})$
shows that the thimble is analytic throughout a region $O(N^{3/2})$
times larger than the BL concentration scale. The overlap is
extremely robust and improves with increasing~$N$.
\end{remark}

\begin{remark}[Lean formalization]
For the initial formalization, we \textbf{axiomatize}
Proposition~\ref{prop:quantum_thimble} (existence of $\Phi$
with bounds). The IFT proof is a realistic future target:
Mathlib has implicit function theorems (both finite-dimensional
and Banach-space versions), and the key hypotheses are available:
smoothness from \texttt{contDiff\_matrix\_det} (proved),
invertibility from $H > 0$ (proved: \texttt{bare\_hessian\_pos}),
and the Hessian formula from \texttt{hessian\_log\_det} (proved).
\end{remark}

\subsubsection{The transfer matrix resolution}
\label{sec:transfer_matrix}

A more robust approach avoids the spacetime phase problem
entirely by reducing the mass gap to a \textbf{spectral gap}
of the transfer matrix.

On a lattice $\La = \La_s \times \{0,\ldots,T-1\}$ (spatial
sites $\times$ time), the partition function factorizes via the
transfer matrix~$\mathcal{T}$:
\begin{equation}
  Z = \Tr(\mathcal{T}^T), \qquad
  \langle\ph(x,t)\ph(0,0)\rangle
    = \frac{\Tr(\mathcal{T}^{T-t}\,\hat\ph(x)\,\mathcal{T}^t\,\hat\ph(0))}
           {\Tr(\mathcal{T}^T)}.
\end{equation}
The mass gap is $m = \log(\mathrm{ev}_0/\mathrm{ev}_1)$, where
$\mathrm{ev}_0 > \mathrm{ev}_1$ are the two largest eigenvalues
of~$\mathcal{T}$. This is a property of $\mathcal{T}$ alone ---
\textbf{independent of~$T$} (the time extent).

The contour shift / thimble analysis is then applied to the
\textbf{spatial} integral that defines~$\mathcal{T}$, which
involves only $|\La_s|$ sites (the spatial volume). For a 2D
theory: $|\La_s| = L/a$ (one spatial dimension), and
$N \gg \sqrt{|\La_s|} = \sqrt{L/a}$ is a condition on $N$
vs.\ the \emph{spatial} volume, not the full spacetime volume.

\textbf{Key advantage}: The spectral gap of $\mathcal{T}$ is an
intrinsic property of the one-step transfer matrix. The phase
problem involves $|\La_s|$ sites (fixed for a given spatial
lattice), not $|\La| = |\La_s| \times T$ (which grows with
the time extent). The mass gap is then:
\begin{enumerate}
\item Proved on a fixed spatial lattice $\La_s$ with
  $N \gg \sqrt{|\La_s|}$ (contour shift for the spatial
  integral in~$\mathcal{T}$).
\item Extended to $|\La_s| \to \infty$ by showing the spectral
  gap is uniform (BL concentration + FK bound, both uniform in
  $|\La_s|$ as argued in Sections~\ref{sec:HJ_solution}--\ref{sec:global_logconcave}).
\item This is exactly the structure of Kupiainen's proof:
  reflection positivity + transfer matrix + cluster expansion
  (replaced here by BL).
\end{enumerate}

\begin{remark}[Three routes to the thermodynamic limit]
\label{rem:three_routes}
The mass gap depends on the ratio of eigenvalues of the
one-step transfer matrix. The sign problem for this matrix
involves only $|\La_s|$ spatial sites --- decoupled from
the temporal extent.

\medskip\noindent
\begin{tabular}{|l|c|l|}
\hline
\textbf{Method} & \textbf{Phase problem} & \textbf{Condition} \\
\hline
Transfer matrix + spatial shift & $|\La_s|$ sites
  & $N \gg \sqrt{L/a}$ \\
Quantum HJ equation & None (exact) & $N > N_0(m_0,\lambda)$ \\
Spacetime contour shift & $|\La|$ sites & $N \gg \sqrt{L^2/a^2}$ \\
\hline
\end{tabular}

\medskip
\begin{enumerate}
\item[(a)] \textbf{Transfer matrix + spatial contour shift}:
  phase problem involves $|\La_s|$ only; $N \gg \sqrt{L/a}$.
  This is the most robust approach because it decouples the
  large-$N$ expansion from the temporal extent of the system.
\item[(b)] \textbf{Quantum HJ equation}: eliminates phase
  exactly; requires second-order PDE existence theory.
  The condition $N > N_0$ depends only on the coupling, not
  on the volume.
\item[(c)] \textbf{Spacetime contour shift with}
  $N \gg \sqrt{|\La|}$: simplest, but couples $N$ to the
  full spacetime volume.
\end{enumerate}
\end{remark}

\section{Solution theory and the mass gap proof}
\label{sec:HJ_solution}

This section establishes the a priori bounds on~$\Phi$ needed for
the mass gap, addresses the technical issues (singularity avoidance,
log-concavity, BL version), and assembles the proof.

\subsection{Solution theory for the HJ equation}

The HJ equation~\eqref{eq:HJ_arctan} is exact but highly nonlinear.
We describe three approaches to its solution and the a priori
bounds needed for the mass gap proof.

\subsubsection{Method of characteristics}

The HJ equation $H(u,\psi) = 0$ with ``Hamiltonian''
$H(u,\psi) = \im f(u+i(v_*+\psi))$ has characteristic equations:
\begin{equation}\label{eq:characteristics}
  \frac{du}{ds} = \frac{\partial H}{\partial\psi}, \qquad
  \frac{d\psi}{ds} = -\frac{\partial H}{\partial u}.
\end{equation}
Since the anti-holomorphic gradient flow~\eqref{eq:flow} preserves
$\im f$, the characteristics~\eqref{eq:characteristics} are
precisely the \textbf{projections of the gradient flow} onto
$(u,\psi)$-space. In other words:

\begin{center}
\emph{Solving the HJ equation and solving the gradient flow are
dual descriptions of the same object.}
\end{center}

The gradient flow constructs the thimble as a stable manifold
(dynamical systems perspective); the HJ equation constructs it
as a Lagrangian graph (symplectic geometry perspective). The
duality means that any estimate proved via the flow automatically
gives an estimate on~$\Phi$, and vice versa.

\subsubsection{A priori bounds on $\Phi$}

For the mass gap proof, we do not need the explicit form of $\Phi$
--- we need three \textbf{a priori bounds}:

\begin{proposition}[Thimble regularity]\label{prop:apriori}
For $N > N_0$ and $\|u\|_\infty \le m_0 N^{1/4}/C$
(inside the convergence disk):
\begin{enumerate}
\item \textbf{Proximity to constant shift}:
  $\displaystyle |\psi(x)| = \Bigl|\frac{\partial\Phi}{\partial u(x)}\Bigr|
    \le \frac{C_1}{\sqrt{N}}\,\|u\|_\infty^2$.
\item \textbf{Spectral gap preservation}:
  $M(\psi) = -\Delta + m_0^2 - 2\psi/\sqrt{N} \;\ge\; m_0^2/2 > 0$.
\item \textbf{Hessian control}:
  $\|\nabla^2\re f|_{\mathcal{J}_0} - H\|_{\mathrm{op}} \le C_2/N$,
  so the Hessian of $\re f$ on the thimble is $H + O(1/N)$.
\end{enumerate}
\end{proposition}

\begin{proof}[Proof sketch]
\textbf{Bound (1)} follows from the convergent Taylor expansion:
$\psi = \nabla\Phi$ starts at order $u^2$ with coefficient
$O(1/\sqrt{N})$ (from the triangle diagram), and higher orders
are controlled by the convergence of the series.

\textbf{Bound (2)} follows from (1): the perturbation
$2|\psi|/\sqrt{N} \le 2C_1\|u\|_\infty^2/N$, which is
$\ll m_0^2$ within the saddle-point region
$\|u\|_\infty \lesssim 1/\sqrt{N}$.

\textbf{Bound (3)}: on the thimble $v = \nabla\Phi(u)$, the
pullback of $\re f$ to $\R^{|\La|}$ (via $u \mapsto u+i\nabla\Phi$)
has Hessian $H + O(1/N)$ corrections from the $\Phi$-dependent
terms in $\re f$ (the cubic and higher terms contribute $O(1/N)$
to the second derivative at the saddle width).
\end{proof}

\textbf{What these bounds give}:
\begin{itemize}
\item Bound (1) + (2): the FK bound applies on the thimble.
  The operator $M(\psi) + iV(u) = -\Delta + m_0^2 + O(1/N) + iV$
  has positive real part $\ge m_0^2/2$, giving
  $|G_\sigma(x,0)| \le C\,e^{-m_0|x|/\sqrt{2}}$ for every
  $u$-configuration on the thimble.
\item Bound (3): Brascamp-Lieb applies on the thimble.
  The measure $e^{\re f}\,du$ has Hessian $\ge \kappa N$
  (with $\kappa > 0$ independent of~$|\La|$), giving
  $\Var(u(x)) \le 1/(\kappa N)$.
\end{itemize}

\subsubsection{Fixed-point formulation}

The HJ equation can be reformulated as a fixed-point problem.
Define the map $\mathcal{T}$ by: given $\psi : \R^{|\La|} \to
\R^{|\La|}$ (a candidate thimble correction), set
\begin{equation}
  (\mathcal{T}\psi)(x)
    = -2\lambda\Bigl[\rho^2 + \frac{N}{2}\,
      \frac{\partial}{\partial u(x)}
      \Tr\arctan\bigl(M(\psi)^{-1/2}\,V(u)\,M(\psi)^{-1/2}\bigr)
    \Bigr] - v_*.
\end{equation}
Then $\psi = \nabla\Phi$ solves the HJ equation if and only if
$\psi$ is a fixed point of~$\mathcal{T}$.

For $N$ sufficiently large: $\mathcal{T}$ is a contraction on the
ball $\|\psi\|_\infty \le m_0^2\sqrt{N}/4$ in an appropriate
function space (the $\arctan$ is Lipschitz with constant $< 1$
when the spectral gap of~$M$ is large). The Banach fixed-point
theorem then gives:
\begin{enumerate}
\item \textbf{Existence}: a unique solution $\psi = \nabla\Phi$
  in the ball.
\item \textbf{Constructive approximation}: the iterates
  $\psi^{(n+1)} = \mathcal{T}\psi^{(n)}$ (starting from
  $\psi^{(0)} = 0$, the constant shift) converge geometrically.
\item \textbf{Quantitative bounds}: $\|\psi - \psi^{(n)}\| \le
  C\,\theta^n$ with contraction rate $\theta < 1$.
\end{enumerate}

This provides an alternative to the Taylor expansion:
a \textbf{non-perturbative}, constructive proof that the thimble
exists and satisfies the a priori bounds of
Proposition~\ref{prop:apriori}.

\subsubsection{What the proof needs vs.\ what we compute}

For the mass gap proof, the HJ equation plays the role of an
\textbf{existence and regularity theorem}, not a computational
tool. Specifically:

\begin{center}
\begin{tabular}{|l|l|}
\hline
\textbf{Step} & \textbf{What is needed from the HJ equation} \\
\hline
Positivity & $\im f = 0$ on the thimble (by construction) \\
BL concentration & Hessian of $\re f \ge \kappa N$ (Bound 3) \\
FK bound & $M(\psi) > 0$ on the thimble (Bound 2) \\
Single-thimble dominance & Analyticity + action gap $\sim N$ \\
\hline
\end{tabular}
\end{center}

None of these require the explicit form of~$\Phi$ --- only
bounds on $\Phi$ and its derivatives, all of which follow from
the convergence of the expansion (or the fixed-point argument).

\subsection{Properties of the integral over $\mathcal{J}_0$}

On $\mathcal{J}_0$, the measure $e^{f(\sigma)}\,d\sigma$ (with the
appropriate orientation) is:
\begin{enumerate}
\item \textbf{Real and positive} (since $\im\,f = 0$
  and $e^{\re\,f} > 0$).
\item \textbf{Log-concave near the saddle}: $\re\,f
  \approx -\tfrac{1}{2}(\sigma,H\sigma)$ (Gaussian, $H > 0$).
\item \textbf{Decaying}: $\re\,f \to -\infty$ along
  $\mathcal{J}_0$ away from the saddle (steepest descent property).
\end{enumerate}

We now address three technical issues in detail.

\subsubsection{Singularity avoidance in the thermodynamic limit}

The singularities of $f$ (where $\det(-\Delta+2i\sigma z) = 0$)
occur at $\sigma = ik^2/(2z)$ --- purely imaginary, at distance
$\ge m_0^2\sqrt{N}/2$ from the saddle. On the thimble,
Bound~(2) of Proposition~\ref{prop:apriori} gives
$M(\psi) \ge m_0^2/2 > 0$, so the thimble \emph{never reaches
these singularities} within the convergence disk.

For extreme fluctuations: BL concentration gives
\begin{equation}
  \Pr\bigl(\|u\|_\infty > t\bigr)
    \le 2|\La|\,e^{-\kappa N t^2/2}.
\end{equation}
The singularity requires $|u| \sim m_0 N^{1/4}$ (the convergence
radius). Substituting $t = m_0 N^{1/4}$:
\begin{equation}
  \Pr\bigl(\text{thimble reaches singularity}\bigr)
    \le 2|\La|\,e^{-\kappa m_0^2 N^{3/2}/2}.
\end{equation}
This is \textbf{superexponentially small in~$N$}, uniformly
in~$|\La|$ (since $e^{-cN^{3/2}}$ beats any polynomial in~$|\La|$
for $N$ large). Singularity avoidance holds in the thermodynamic
limit.

\subsubsection{Global vs.\ local log-concavity}
\label{sec:global_logconcave}

The Hessian of $\re f$ on the thimble is $H + O(1/N)$
(Bound~3) \emph{near the saddle}. Far from the saddle, the
$\Tr\log$ term can introduce non-convexity. The question: does
BL require \emph{global} log-concavity?

\textbf{Answer}: For our application, \emph{local} log-concavity
within the Gaussian peak suffices, for two reasons:

\begin{enumerate}
\item \textbf{The Gaussian dominates the tails}. The bare term
  $-\|u\|^2/(4\lambda)$ in $\re f$ provides Gaussian decay.
  The $\Tr\log$ correction to the Hessian is bounded by
  $\|(\partial^2/\partial u^2)\Tr\log M\| \le C/m_0^4$
  (from the resolvent bound $\|M^{-1}\| \le 2/m_0^2$).
  The Gaussian Hessian $1/(2\lambda)$ dominates for
  $\lambda < m_0^4/(2C)$ --- a condition on the coupling,
  not the volume.

\item \textbf{BL with a cutoff}. One can apply BL to the
  measure restricted to the log-concave region
  $\|u\|_\infty \le R$ (with $R \sim m_0 N^{1/4}/C$, the
  convergence radius), and bound the contribution from
  $\|u\|_\infty > R$ by the Gaussian tail. Since the tail
  probability is $\le 2|\La|e^{-\kappa NR^2/2}
  \sim e^{-cN^{3/2}}$, the correction is negligible.
\end{enumerate}

In fact, for $\lambda$ small enough (weak coupling), $\re f$ on
the thimble IS globally strictly convex: the Gaussian $1/(2\lambda)$
term dominates the $\Tr\log$ curvature everywhere. The threshold
is $\lambda < c\,m_0^4$, which is satisfied in the weak-coupling
regime relevant to the continuum limit.

\subsubsection{Brascamp-Lieb: which version?}

The measure $e^{\re f}\,du$ on the thimble is a probability
measure on~$\R^{|\La|}$ (after normalization). For BL, we need:

\begin{itemize}
\item \textbf{Standard BL} (Brascamp-Lieb 1976~\cite{BL76}):
  applies to $e^{-V(u)}\,du$ with $V$ strictly convex
  ($\nabla^2 V \ge \kappa > 0$). This works on a finite
  lattice if $V = -\re f$ is globally strictly convex ---
  which holds for $\lambda < c\,m_0^4$ (see above).
  \textbf{No extension needed} in this regime.

\item \textbf{Conlon-Dabkowski extension}~\cite{CD24}:
  applies to lattice field models with uniformly convex
  \emph{gradient} action (weaker than strict convexity of~$V$).
  This would be needed if $\lambda$ is not small, i.e., the
  non-convex $\Tr\log$ tails are significant. Their result
  is stated for lattice models in $d \ge 2$ and gives
  near-Gaussian correlation bounds.

\item \textbf{BL from markov-semigroups} (our Lean library):
  the \texttt{brascampLieb\_poincare} theorem gives
  $\Var(f) \le (1/(\kappa N))\,\E[\|\nabla f\|^2]$ for
  $(\kappa N)$-log-concave measures. This is the version we
  would formalize. It requires \texttt{LogConcaveMeasure}
  with convexity parameter~$\kappa N$, which we construct
  from the Hessian bound~$H \ge \kappa N$.
\end{itemize}

\subsubsection{Uniformity and the infinite-volume limit}
\label{sec:uniformity}

The mass gap is the spectral gap of the transfer matrix
$\mathcal{T}$ (Section~\ref{sec:transfer_matrix}). The contour
shift / BL analysis is applied to the \textbf{spatial} integral
defining~$\mathcal{T}$, which involves $|\La_s|$ sites.

The proof requires $N \gg \sqrt{|\La_s|}$ to control the
residual phase. For a fixed spatial lattice (e.g.,
$|\La_s| = L/a$ in 2d), this is a \textbf{fixed} condition
$N > N_0(L/a, m_0, \lambda)$ independent of the time extent.

\textbf{Taking $|\La_s| \to \infty$ (spatial thermodynamic limit)}:
The spectral gap of $\mathcal{T}$ must be shown uniform in~$|\La_s|$.
This is a stronger statement requiring either:
\begin{itemize}
\item[(a)] The quantum HJ equation~\eqref{eq:quantum_HJ}
  (eliminating the phase), or
\item[(b)] A direct proof that the spectral gap of $\mathcal{T}$
  is controlled by the per-site BL + FK bounds (which are
  volume-independent), without passing through the phase.
\end{itemize}
Kupiainen's proof uses option~(b) via a cluster expansion.
The BL approach should give the same result: the FK bound is
per-configuration (volume-independent), and BL concentration
holds with $\kappa$ independent of~$|\La_s|$ (the Hessian
$\hat{H}(k) = 1/(2\lambda) + 2B(k)$ has a spectral gap
converging to a finite limit as $|\La_s| \to \infty$).

The remaining issue is the \textbf{phase control} in the spatial
integral. For the transfer matrix, the relevant phase is
$O(|\La_s|/N^2)$, which vanishes for $N \gg \sqrt{|\La_s|}$.
In the continuum limit $a \to 0$ (so $|\La_s| = L/a \to \infty$),
this requires $N \gg \sqrt{L/a}$ --- coupling $N$ to the
UV cutoff.

\begin{remark}[The continuum limit is weak coupling]
\label{rmk:weak_coupling}
The global log-concavity condition $\lambda < c\,m_0^4$ might
appear restrictive, since $m_0 \sim \La\,e^{-c/g^2}$ is
exponentially small. But the \textbf{continuum limit is
automatically in the weak-coupling regime}: the renormalized
coupling runs as $g^2(a) \sim 1/\log(\La/m_0) \to 0$ as
$a \to 0$. The bare coupling $\lambda \sim g^2/N$, while
$m_0^4 \sim \La^4\,e^{-4c/g^2}$, so:
\begin{equation}
  \frac{\lambda}{m_0^4}
    \;\sim\; \frac{g^2}{N\La^4\,e^{-4c/g^2}}
    \;\to\; 0
    \qquad\text{as } a \to 0.
\end{equation}
The condition $\lambda < c\,m_0^4$ is satisfied with
\textbf{enormous margin} in the continuum limit. Global
log-concavity is a consequence of asymptotic freedom, not an
extra assumption.

The Conlon-Dabkowski extension~\cite{CD24} would be needed for
strong coupling (fixed lattice spacing, $\lambda$ not small) or
for $N$ close to the threshold~$N_0$. Neither is relevant for
the initial proof, which targets the continuum limit at
large~$N$.
\end{remark}

\subsection{The flowed manifold alternative}

Instead of characterizing the exact thimble $\mathcal{J}_0$
(a global object determined by the full analytic structure of~$f$),
one can work with the \textbf{flowed manifold}:
\begin{equation}
  \mathcal{M}_T = \bigl\{\sigma(T) : \sigma(0) \in \R^{|\La|},\;
    \dot\sigma = \overline{\nabla f(\sigma)}\bigr\},
\end{equation}
obtained by flowing the real domain along~\eqref{eq:flow} for
time~$T$. Properties:
\begin{itemize}
\item $\mathcal{M}_0 = \R^{|\La|}$ (the original domain).
\item $\mathcal{M}_T$ is diffeomorphic to $\R^{|\La|}$ for all
  finite~$T$ (no topology change).
\item $\mathcal{M}_T \to \bigcup_{\sigma_*}\mathcal{J}_{\sigma_*}$
  as $T \to \infty$.
\item The residual sign problem on $\mathcal{M}_T$ is bounded by
  $e^{-c(T)N}$ for a function $c(T)$ increasing in~$T$.
\end{itemize}
\textbf{Advantage}: No need to identify critical points or compute
intersection numbers. The flow~\eqref{eq:flow} is a well-posed ODE
system. For a rigorous proof, take $T$ large enough that the residual
oscillations are smaller than the mass gap signal.

This approach was developed computationally by
Fukuma et al.~\cite{Fukuma20,Fukuma23} (``Worldvolume HMC''),
now extended to gauge theories~\cite{Fukuma26} and the Hubbard
model~\cite{Fukuma25b}.

\subsection{Proof plan (the quantum thimble argument)}

The quantum thimble (Section~\ref{sec:quantum_HJ}) gives a
remarkably simple proof of the mass gap:

\begin{enumerate}
\item \textbf{Contour deformation}: Deform the $\sigma$-integral
  from $\R^{|\La|}$ to the quantum thimble, where the total
  integrand $e^f \cdot \det J$ is \textbf{real and positive}
  (quantum HJ eliminates both action and Jacobian phases).
  The two-point function becomes:
  \begin{equation}
    \langle\ph(x)\ph(0)\rangle
      = \frac{1}{Z}\int G_\sigma(x,0)\,e^{-V_{\rm eff}(u)}\,du
  \end{equation}
  where $e^{-V_{\rm eff}}$ is a positive measure on~$\R^{|\La|}$.

\item \textbf{FK bound (uniform in $u$)}: On the quantum thimble,
  the $\ph$-operator is $-\Delta + m_0^2 + 2iuz$ with $u \in \R$.
  The FK/diamagnetic bound gives
  $|G_\sigma(x,0)| \le G_{m_0}(x,0)$ for \textbf{every}
  $u$-configuration. This is the operator monotonicity
  $|(\text{M}+i\text{V})^{-1}| \le \text{M}^{-1}$.

\item \textbf{Trivial averaging}: Since $e^{-V_{\rm eff}}$ is a
  \textbf{positive} measure and the FK bound is uniform in~$u$:
  \begin{equation}
    |\langle\ph(x)\ph(0)\rangle|
      \le \frac{1}{Z}\int |G_\sigma|\,e^{-V_{\rm eff}}\,du
      \le G_{m_0}(x,0)\cdot\frac{Z}{Z}
      = G_{m_0}(x,0).
  \end{equation}
  The ratio $Z/Z = 1$ because the measure is positive ---
  \textbf{no sign problem, no average sign factor}.

\item \textbf{Green's function decay}: $G_{m_0}(x,0) =
  (-\Delta+m_0^2)^{-1}(x,0) \le (2/m_0^2)\,e^{-m_0|x|}$
  (Combes-Thomas conjugation estimate for finite-range operators;
  see \S\ref{sec:combes_thomas}).

\item \textbf{Mass gap}: Chaining steps 1--4:
  $|\langle\ph(x)\ph(0)\rangle| \le (2/m_0^2)\,e^{-m_0|x|}$
  with $m_0 > 0$ from the gap equation, independent of~$|\La|$.
\end{enumerate}

\begin{remark}[The role of BL in the mass gap proof]
The BL variance bound $\Var(u(x)) \le 1/(\kappa N)$ is not
used in the mass gap chain (steps 2--4): the FK bound is
uniform in~$u$, so the average is controlled by the triangle
inequality on the positive measure.

However, BL enters \emph{indirectly}: the quantum thimble
existence (step~1) requires the Hessian spectral gap $\kappa > 0$
via the implicit function theorem. The Hessian bound IS the
BL convexity parameter. So the logical dependencies are:
\begin{center}
Hessian $\kappa > 0$ $\;\xrightarrow{\text{IFT}}\;$
quantum thimble exists $\;\xrightarrow{\text{positive measure}}\;$
FK bound averages trivially $\;\to\;$ mass gap.
\end{center}
The Hessian bound $\hat{H}(k) = 1/(2\lambda) + 2B(k) \ge \kappa > 0$
is uniform in~$|\La|$ (the spectral gap converges to a finite
limit). So the quantum thimble exists for all finite~$|\La|$
simultaneously, and the mass gap is uniform in volume.
\end{remark}

\begin{remark}[Volume independence]
The constants $m_0$ and $2/m_0^2$ depend on the gap equation
(coupling constants), \textbf{not on~$|\La|$}. The mass gap
persists in the thermodynamic limit because:
\begin{itemize}
\item FK bound: per-configuration, no $|\La|$ dependence.
\item Positive measure: quantum HJ, ratio $Z/Z = 1$.
\item Green's decay: $m_0$ from gap equation, independent of~$|\La|$.
\end{itemize}

For a rigorous proof, one can work with the \emph{exact} constant
shift (Cauchy's theorem for a flat contour is elementary) and treat
the difference between the constant shift and the exact thimble
as a perturbation. The residual $\im\,f$ on the constant-shift
contour is $O(u^3/\sqrt{N})$ per site --- small within the
saddle-point width but requiring control in the tails.
\end{remark}

\begin{remark}[Decomposition of the correlator bound --- formalization status]
\label{rem:decomposition}
The correlator bound (Step~2 of the proof plan) decomposes into
four sub-steps, of which only one is research-level:

\begin{enumerate}
\item[\textbf{S1.}] \textbf{HS representation}: $\langle\ph\ph\rangle_c
  = (1/Z)\int G_\sigma \cdot w(\sigma)\,d\sigma$.
  \textbf{Proved.} The partition function version (\texttt{hs\_partition\_complex})
  is fully proved; the correlator version follows from the Gaussian two-point
  formula (available in gaussian-field).

\item[\textbf{S2.}] \textbf{Cauchy contour shift}: $\int_\R \to \int_{\text{thimble}}$.
  \textbf{Proved ingredients.} Single-variable \texttt{vertical\_contour\_shift}
  is fully proved from Mathlib's Cauchy integral theorem. Multi-variable
  follows from Fubini plus Gaussian decay of the bare weight.

\item[\textbf{S3.}] \textbf{Thimble positivity}: $e^f \cdot \det J > 0$ on the thimble.
  \textbf{OPEN} (\texttt{quantum\_thimble\_exists}).
  This requires the implicit function theorem for the quantum Hamilton-Jacobi
  equation (Section~\ref{sec:HJ}) plus the Brascamp-Lieb variance bound.
  This is the only research-level gap.

\item[\textbf{S4.}] \textbf{Triangle + FK averaging}: on a positive measure,
  $|\int G \cdot \mu| \le \int |G| \cdot \mu \le M^{-1}$.
  \textbf{Trivial} (standard measure theory + diamagnetic \S\ref{sec:diamagnetic}).
\end{enumerate}

The \texttt{correlator\_le\_thimble\_avg} axiom in Lean packages S1--S4.
All content except S3 is proved or trivially provable. The mass gap proof
depends on exactly two research-level results: (1)~thimble positivity
from the quantum HJ equation, and (2)~the diamagnetic inequality.
\end{remark}

\section{The diamagnetic inequality}\label{sec:diamagnetic}

The FK bound $|G_\sigma(x,0)| \le G_{m_0}(x,0)$ used in Step~3
of the proof plan (\S10.2) relies on the \textbf{diamagnetic
inequality} for the resolvent of $M + iV$, where $M = -\Delta +
m_0^2$ is positive definite and $V = 2uz$ is a real diagonal
matrix. This section presents the semigroup proof, following
Simon~\cite{Simon1979} Ch.~22 and Reed--Simon~IV \S X.4.

\subsection{Statement}

\begin{theorem}[Diamagnetic inequality for finite matrices]
\label{thm:diamagnetic}
Let $M$ be a real symmetric positive definite $n \times n$ matrix
with nonpositive off-diagonal entries (a Stieltjes matrix), and let
$V$ be a real diagonal matrix. Then for all $x, y$:
\begin{equation}\label{eq:diamagnetic}
  |(M + iV)^{-1}(x,y)| \le M^{-1}(x,y).
\end{equation}
In particular, $M^{-1}(x,y) \ge 0$ for all $x,y$.
\end{theorem}

In our application, $M = -\Delta + m_0^2$ (the shifted lattice
Laplacian) and $V = \mathrm{diag}(2u(x)/\sqrt{N})$, so the
hypotheses are satisfied: $M$ has off-diagonal entries $-1$ (from
the nearest-neighbor Laplacian) and diagonal entries $2d + m_0^2
\ge m_0^2 > 0$, giving positive definiteness by Gershgorin.

\subsection{The semigroup proof}

The proof decomposes into five steps.

\medskip\noindent\textbf{Step 1: Laplace transform representation.}
For any positive definite matrix~$A$,
\begin{equation}\label{eq:laplace-resolvent}
  A^{-1} = \int_0^\infty e^{-tA}\,dt.
\end{equation}
This follows from spectral decomposition: $A = \sum_k \lambda_k
|e_k\rangle\langle e_k|$ with all $\lambda_k > 0$, so
$\int_0^\infty e^{-t\lambda_k}\,dt = 1/\lambda_k$. Similarly,
$(M+iV)^{-1} = \int_0^\infty e^{-t(M+iV)}\,dt$ since the
eigenvalues of $M + iV$ have positive real part (from $M > 0$).

\medskip\noindent\textbf{Step 2: Lie--Trotter product formula.}
For finite matrices $A$ and $B$,
\begin{equation}\label{eq:trotter}
  e^{A+B} = \lim_{n\to\infty} \bigl(e^{A/n}\,e^{B/n}\bigr)^n.
\end{equation}
This is the Lie product formula; for matrices it follows from
$e^{A/n}e^{B/n} = e^{(A+B)/n + O(1/n^2)}$ (Baker--Campbell--Hausdorff)
and $(I + X/n)^n \to e^X$. Applied to $A = -tM$ and $B = -itV$:
\begin{equation}
  e^{-t(M+iV)} = \lim_{n\to\infty}
    \bigl(e^{-tM/n}\,e^{-itV/n}\bigr)^n.
\end{equation}

\medskip\noindent\textbf{Step 3: Phase bound.}
Since $V = \mathrm{diag}(v_1, \ldots, v_n)$ is diagonal,
\begin{equation}
  e^{-itV/n}(x,y) = \delta_{xy}\,e^{-itv_x/n}
\end{equation}
has $|e^{-itV/n}(x,y)| = \delta_{xy}$. In particular, for any
matrix~$P$,
\begin{equation}\label{eq:phase-bound}
  |(P\,e^{-itV/n})(x,y)| = |P(x,y)\,e^{-itv_y/n}| = |P(x,y)|.
\end{equation}

\medskip\noindent\textbf{Step 4: Heat kernel positivity.}
Since $M$ has nonpositive off-diagonal entries (it is a \emph{Z-matrix}
or equivalently $-M$ is a \emph{Metzler matrix}), the matrix
exponential $e^{-tM}$ has nonneg entries for all $t \ge 0$:
\begin{equation}\label{eq:heat-nonneg}
  e^{-tM}(x,y) \ge 0 \quad\text{for all } x, y, t \ge 0.
\end{equation}

\begin{proof}[Proof of~\eqref{eq:heat-nonneg}]
Write $-M = D + R$ where $D$ is the diagonal part and $R \ge 0$
entrywise (since $M$ has nonpositive off-diagonal). By the
Euler approximation:
\begin{equation}
  e^{-tM} = \lim_{n\to\infty} \Bigl(I - \frac{tM}{n}\Bigr)^n
  = \lim_{n\to\infty} \Bigl(I + \frac{t(D+R)}{n}\Bigr)^n.
\end{equation}
For $n > t\|M\|$, each factor $I + t(D+R)/n$ has nonneg entries
(diagonal entries $\ge 1 - t\|D\|_\infty/n > 0$, off-diagonal
entries $= tR_{xy}/n \ge 0$). The product of entrywise-nonneg
matrices is entrywise-nonneg, so $e^{-tM}(x,y) \ge 0$.
\end{proof}

\begin{remark}[Metzler shift proof (used in Lean)]
An alternative proof avoids the Euler limit: write $-M = (-M+\alpha I)
- \alpha I$ where $\alpha \ge \max_i M_{ii}$. The matrix $C = -M + \alpha I
\ge 0$ entrywise (off-diagonal nonneg by Z-matrix, diagonal shifted to nonneg).
Then $e^{-tM} = e^{-t\alpha} \cdot e^{tC}$ (since $\alpha I$ commutes with
everything). Now $e^{tC} = \sum_{k\ge 0} (tC)^k/k! \ge 0$ (power series of
nonneg matrix with nonneg coefficients) and $e^{-t\alpha} > 0$.
This is fully proved in Lean without the Euler convergence step.
\end{remark}

\medskip\noindent\textbf{Step 5: Assembly.}
Each Trotter factor satisfies:
\begin{equation}
  |(e^{-tM/n}\,e^{-itV/n})(x,y)|
  \stackrel{\eqref{eq:phase-bound}}{=} |e^{-tM/n}(x,y)|
  \stackrel{\eqref{eq:heat-nonneg}}{=} e^{-tM/n}(x,y).
\end{equation}
For the $n$-fold product, the entry-wise triangle inequality gives:
\begin{align}
  |\bigl(e^{-tM/n}e^{-itV/n}\bigr)^n(x,y)|
  &\le \sum_{z_1,\ldots,z_{n-1}}
    \prod_{k=1}^n |(e^{-tM/n}e^{-itV/n})(z_{k-1},z_k)| \\
  &= \sum_{z_1,\ldots,z_{n-1}}
    \prod_{k=1}^n e^{-tM/n}(z_{k-1},z_k) \\
  &= (e^{-tM/n})^n(x,y) = e^{-tM}(x,y).
\end{align}
Taking $n \to \infty$ (Trotter) and integrating over $t$
(Laplace transform):
\begin{equation}
  |(M+iV)^{-1}(x,y)|
  = \Bigl|\int_0^\infty e^{-t(M+iV)}(x,y)\,dt\Bigr|
  \le \int_0^\infty |e^{-t(M+iV)}(x,y)|\,dt
  \le \int_0^\infty e^{-tM}(x,y)\,dt
  = M^{-1}(x,y). \qed
\end{equation}

\begin{remark}[M-matrix inverse nonnegativity]
As a corollary, $M^{-1}(x,y) \ge 0$: from~\eqref{eq:laplace-resolvent}
and~\eqref{eq:heat-nonneg}, $M^{-1} = \int_0^\infty e^{-tM}\,dt
\ge 0$ entrywise. This is a standard result in M-matrix theory.
\end{remark}

\begin{remark}[The Lean formalization]
The Lean file \texttt{DiagmagneticInequality.lean} (413 lines)
decomposes this proof into 6 axioms corresponding to Steps~1--5
above, plus the M-matrix corollary. The phase bound (Step~3) and
the entry-wise triangle inequality are fully proved; the remaining
axioms (Trotter, heat kernel positivity, Laplace transform) use
standard Mathlib tools (matrix exponential, spectral theorem).
\end{remark}

\section{Green's function decay: the Combes-Thomas estimate}
\label{sec:combes_thomas}

The exponential decay $G_{m_0}(x,0) \le (2/m_0^2)\,e^{-m_0|x|}$
used in Step~4 of the proof plan follows from a general result:

\begin{theorem}[Combes-Thomas estimate]\label{thm:combes_thomas}
Let $M$ be a real symmetric positive definite matrix on a finite
set~$\La$ with spectral gap~$\gamma > 0$ and finite range~$R$
(i.e., $M(x,y) = 0$ when $d(x,y) > R$). Then
\begin{equation}
  |M^{-1}(x,y)| \le C\,e^{-\alpha\,d(x,y)}
\end{equation}
where $\alpha = R^{-1}\log(1 + \gamma/(2\|M\|))$ and
$C = 2|\La|/\gamma$.
\end{theorem}

\begin{proof}
Fix $y_0 \in \La$ and conjugate $M$ by the diagonal matrix
$D = \mathrm{diag}(e^{\alpha\,d(x,y_0)})$:
\[
  M_\alpha = D\,M\,D^{-1}.
\]
Since $M$ has range~$R$, the perturbation satisfies
$|M_\alpha(x,z) - M(x,z)| \le (e^{\alpha R}-1)\,|M(x,z)|$
(using the triangle inequality $|d(x,y_0)-d(z,y_0)| \le d(x,z) \le R$).
Thus $\|M_\alpha - M\| \le (e^{\alpha R}-1)\|M\|$.

Choosing $\alpha$ so that $(e^{\alpha R}-1)\|M\| = \gamma/2$ gives
$M_\alpha$ a spectral gap $\ge \gamma/2$. Since
$M^{-1}(x,y_0) = e^{-\alpha\,d(x,y_0)}\,M_\alpha^{-1}(x,y_0)$,
the bound follows from $|M_\alpha^{-1}(x,y_0)| \le \|M_\alpha^{-1}\|
\le |\La|/(\gamma/2) = C$.
\end{proof}

For the mass gap application: $M = -\Delta + m_0^2$ is the
nearest-neighbor Laplacian on the $d$-dimensional torus, with
spectral gap~$m_0^2$, range~$R = 1$ (nearest-neighbor), and
$\|M\| \le 2d + m_0^2$ (Gershgorin). The decay rate is
$\alpha = \log(1 + m_0^2/(4d + 2m_0^2))$, which matches
the Fourier-computed rate for $d = 1$.

\begin{remark}[Comparison with Fourier analysis]
For $d=1$ (the 1D torus $\Z/L\Z$), the Green's function has
the explicit formula
$G(n) = (r_-^n + r_-^{L-n})/((r_+ - r_-)(1 - r_-^L))$
where $r_\pm = ((2+m^2) \pm \sqrt{m^4+4m^2})/2$ are the
characteristic roots. The decay rate $-\log r_-$ matches the
Combes-Thomas rate for $d=1$. The sharp constant is $2/m^2$
(not $1/m^2$, which fails for $L=2$).

For $d \ge 2$, the Fourier approach is harder (the 2D Green's
function does not factor as a product of 1D Green's functions),
while Combes-Thomas applies unchanged.
\end{remark}

\begin{remark}[The Lean formalization]
The Combes-Thomas estimate is fully proved (0~sorries) in
\texttt{gaussian-field/Lattice/CombesThomas.lean} (760 lines).
The 1D explicit formula is fully proved in
\texttt{pphi2N/Thimble/GreenDecay.lean} via operator verification
(both the Fourier Green's function and the explicit formula satisfy
$(-\Delta+m^2)G = \delta_0$) and positive definiteness injectivity
(both solutions must agree).
\end{remark}

\section{Other strategies for comparison}

Section~\ref{sec:thimble} describes what we call \textbf{Strategy~D}
(Lefschetz thimble). For context, we compare with three alternatives.

\subsection{Strategy A: NLSM (Kupiainen)}

Work with the NLSM directly. The constraint gives real mass $m_0^2$
for free. The FK bound gives $|G| \le G_{m_0}$ (single-$\sigma$
decay). Cluster expansion + RP control the $\sigma$-average.

\textbf{Advantage}: Real mass in the operator, clean FK bound.\\
\textbf{Disadvantage}: Cluster expansion needed (technically heavy).\\
\textbf{Threshold}: $N > C\,m_0^{-a_2}$ with $a_2 = O(d^2)$
(Kupiainen Thm.~1).

\subsection{Strategy B: LSM saddle-point (real axis)}

Work with the LSM on the \emph{real} $\sigma$-axis. The saddle-point
Gaussian is log-concave (real, positive), but the full integrand has
oscillatory corrections away from the saddle.

\textbf{Idea}: Apply BL to the saddle-point Gaussian and control the
oscillatory corrections as a perturbation.

\textbf{Challenge}: The mass gap comes entirely from the $\sigma$-average
(no single-$\sigma$ decay). Need to show the oscillatory corrections
don't destroy the saddle-point mass.

\textbf{Advantage}: Avoids cluster expansion (BL for the Gaussian).\\
\textbf{Disadvantage}: No FK bound; the sign problem is addressed
by perturbative control, not eliminated.\\
\textbf{Comparison with D}: Strategy~D \emph{eliminates} the
oscillations by deforming the contour; Strategy~B tries to
\emph{control} them on the real axis.

\subsection{Strategy C: LSM $\to$ NLSM equivalence}

Prove that the LSM and NLSM have the same IR physics (they differ
only at the $\sigma$-mass scale $M_\sigma \sim \sqrt{\lambda}\,v$),
then use Strategy~A.

\textbf{Advantage}: Gets the NLSM mass gap (with the real mass trick)
from the LSM formulation.\\
\textbf{Disadvantage}: Proving IR equivalence rigorously is itself
a major result.

\subsection{Strategy comparison}

\begin{center}
\begin{tabular}{|l|c|c|c|c|}
\hline
& A (NLSM) & B (real axis) & C (equivalence) & D (thimble) \\
\hline
Sign problem & avoided & controlled & avoided & reduced/eliminated \\
Real mass $m_0^2$ & from constraint & none & from constraint & from contour shift \\
FK bound & yes & no & yes & yes \\
BL needed & no & yes & no & yes \\
Cluster exp.\ & yes & no & yes & no \\
Proved & yes (K80) & no & no & no \\
\hline
\end{tabular}
\end{center}

Strategy~D with the constant shift (Section~\ref{sec:constant_shift})
achieves the same operator structure as Strategy~A ($-\Delta+m_0^2$
in the real part) without the constraint. The key insight: \emph{the
thimble does for the LSM what the constraint does for the NLSM}.
The residual sign problem on the constant-shift contour is
$O(|\La|/N^2)$, controlled perturbatively for $N \gg \sqrt{|\La|}$.
This avoids both the cluster expansion (Strategy~A) and the
uncontrolled oscillations (Strategy~B).

\section{The $N = 2$ obstruction}

For $N = 2$: the BKT transition occurs at finite $\lambda$
(Dario-Garban~\cite{DG25}), driven by vortices ($\pi_1(S^1)=\Z$).
Any proof must exclude $N = 2$. For $N \ge 3$: $\pi_1(S^{N-1})=0$,
no BKT.

\section{Current status}

\subsection{Main theorem (in Lean)}

\textbf{42 files, 17 axioms, 0 sorries.}

\begin{theorem}[\texttt{ON\_LSM\_hasCorrelationDecay}]
For the O($N$) LSM interacting measure $\mu$ on a finite lattice
with gap equation mass $m_0^2 > 0$:
\begin{equation}
  \texttt{HasCorrelationDecay}\;\mu\;\mathrm{dist}
\end{equation}
i.e., $\exists\,m, C > 0$ such that
$|\langle\ph^i(x)\ph^i(y)\rangle_c| \le C\,e^{-m\,\mathrm{dist}(x,y)}$
for all components~$i$ and sites $x, y$.
\end{theorem}

The proof chains two axioms:
\begin{enumerate}
\item \texttt{correlator\_le\_thimble\_avg}: the connected correlator
  of the O($N$) LSM measure is bounded by the thimble average of
  $|G_\sigma|$, which is $\le M^{-1}(x,y)$ by the FK bound
  (bundled in \texttt{ThimbleIntegralData}).
\item \texttt{green\_exponential\_decay}: $M^{-1}(x,y) \le
  (2/m_0^2)\,e^{-m_0\,\mathrm{dist}(x,y)}$.
\end{enumerate}

\subsection{What is proved (0 sorries)}

\begin{itemize}
\item \textbf{HS identity}: from Mathlib's
  \texttt{fourierIntegral\_gaussian}.
\item \textbf{HS equivalence} (one-site and multi-site):
  $\int d\sigma\,e^{-S_{\rm HS}} = c\cdot e^{-S_{\rm original}}$
  --- via \texttt{push\_cast; ring} and Fubini.
\item \textbf{Contour shift}: \texttt{vertical\_contour\_shift}
  proved from Mathlib's Cauchy integral theorem (rectangle integral
  + limits + decay).
\item \textbf{Gap equation algebra}: $-2v_* z = m_0^2$, $v_* < 0$,
  singularity distance, physical mass $m_0 > 0$.
\item \textbf{Shifted operator}: $M = -\Delta + m_0^2$ has spectral
  gap $\ge m_0^2$ (from \texttt{psd\_add\_scalar\_bound}).
\item \textbf{Phase cancellation}: $r_1 e^{i\theta} \cdot r_2 e^{-i\theta}$
  is real and positive (polar form).
\item \textbf{1D diamagnetic inequality}: $a \le \|a + bi\|$.
\item \textbf{Lattice Green's function}: $\|G(n)\| \le 1/m^2$
  (17 theorems from Mathlib Fourier API on \texttt{ZMod}),
  plus exponential decay $\|G(n)\| \le (2/m^2) r_-^{d(n)}$
  (from the explicit formula axiom via Vieta's identity).
\item \textbf{det is $C^\infty$}: proved with linftyOp norm
  via Leibniz formula + entry bound (was axiom).
\item \textbf{Main theorem}: \texttt{ON\_LSM\_hasCorrelationDecay}
  (from 2 axioms: \texttt{correlator\_le\_thimble\_avg} +
  \texttt{green\_exponential\_decay}).
\item \textbf{Continuum limit}: OS0--OS2 for the O($N$) LSM on $T^2_L$.
\end{itemize}

\subsection{What is axiomatized (17 axioms)}

\begin{center}
\begin{tabular}{|l|l|c|}
\hline
\textbf{Axiom} & \textbf{Content} & \textbf{Difficulty} \\
\hline
\multicolumn{3}{|c|}{\emph{Mass gap chain (3 axioms)}} \\
\hline
\texttt{correlator\_le\_thimble\_avg} & HS + Cauchy + triangle on thimble
  & medium \\
\texttt{green\_exponential\_decay} & $M^{-1} \le Ce^{-m_0|x|}$ (concrete)
  & medium \\
\texttt{greenFunction\_explicit\_formula} & $G(n)$ closed form on $\Z/L\Z$
  & medium \\
\hline
\multicolumn{3}{|c|}{\emph{Diamagnetic inequality (\S\ref{sec:diamagnetic}, 7 axioms)}} \\
\hline
\texttt{resolvent\_complex\_bound} & $|(M{+}iV)^{-1}| \le M^{-1}$
  & medium \\
\texttt{heat\_kernel\_entrywise\_nonneg} & $e^{-tM} \ge 0$ (Z-matrix)
  & medium \\
\texttt{trotter\_product\_matrix} & Lie--Trotter for matrices
  & medium \\
\texttt{diamagnetic\_inequality} & $|e^{-t(M{+}iV)}| \le e^{-tM}$
  & medium \\
\texttt{laplace\_transform\_inverse} & $M^{-1} = \int e^{-tM}\,dt$
  & medium \\
\texttt{laplace\_transform\_inverse\_complex} & $(M{+}iV)^{-1} = \int\cdots$
  & medium \\
\texttt{m\_matrix\_inverse\_nonneg} & $M^{-1} \ge 0$ (M-matrix)
  & easy \\
\hline
\multicolumn{3}{|c|}{\emph{Quantum thimble (1 axiom)}} \\
\hline
\texttt{quantum\_thimble\_exists} & QHJ solution + BL variance
  & hard \\
\hline
\multicolumn{3}{|c|}{\emph{Continuum limit (4 axioms, ports)}} \\
\hline
4 embedding/limit axioms & from pphi2/gaussian-field
  & routine \\
\hline
\multicolumn{3}{|c|}{\emph{Matrix calculus (2 axioms)}} \\
\hline
\texttt{fderiv\_log\_det} & Jacobi's formula
  & medium \\
\texttt{hessian\_log\_det} & Hessian of $\log\det$
  & medium \\
\hline
\end{tabular}
\end{center}

\subsection{Open questions}

\begin{enumerate}
\item \textbf{Singularity avoidance}: Can we prove $\mathcal{J}_0$
  avoids $\det(-\Delta+2i\sigma z) = 0$? The zeros are purely
  imaginary ($\sigma = ik^2/(2z)$), at distance $\ge m_0^2/(2z)
  = m_0^2\sqrt{N}/2$ from the saddle. For large~$N$ this should
  be straightforward.
\item \textbf{Global log-concavity}: Is $\re\,f$ concave on all
  of $\mathcal{J}_0$, or only near the saddle? The $\Tr\log$ term
  introduces non-convexity far from the saddle. For BL: we need
  convexity within the $O(1/\sqrt{N})$-width peak, which holds by
  the Hessian computation.
\item \textbf{Intersection numbers}: Can we verify $n_0 = 1$?
  For the dominant saddle on the real axis with $f''$ positive
  definite, homotopy arguments should give $n_0 = 1$.
\item \textbf{Threshold $N_0$}: What is the minimal $N$ for
  single-thimble dominance? Kupiainen's threshold for Strategy~A
  is $N > C\,m_0^{-25}$ (very large). Can Strategy~D improve this?
\item \textbf{$N = 1$ test case}: In $d = 0$ (one site),
  $\mathcal{J}_0$ is a curve in~$\C$ that can be computed
  explicitly. Does the thimble analysis reproduce the known
  P($\ph$)$_2$ mass gap?
\item \textbf{Flowed manifold vs.\ exact thimble}: Is the flowed
  manifold $\mathcal{M}_T$ for finite~$T$ sufficient for the
  proof, avoiding the need to characterize $\mathcal{J}_0$ globally?
\end{enumerate}

\subsection{Gap between this document and the Lean formalization}

The Lean development mirrors the top-level proof architecture
(Sections~\ref{sec:thimble}--\ref{sec:HJ_solution}) with significant
proved content:
\begin{itemize}
\item The \textbf{main theorem} \texttt{ON\_LSM\_hasCorrelationDecay}
  is proved from 2 axioms. The correlator bound is decomposed into
  \texttt{ThimbleIntegralData} (bundling the thimble average with
  the FK bound) plus the Green's function decay.
\item The \textbf{contour shift} (\texttt{vertical\_contour\_shift})
  is fully proved from Mathlib's Cauchy integral theorem.
\item The \textbf{HS equivalence} (both one-site and multi-site)
  is fully proved.
\item The \textbf{Green's function} (\texttt{GreenDecay.lean}, 460+
  lines) has 17 proved theorems including the sharp exponential
  decay $\|G(n)\| \le (2/m^2) r_-^{d(n)}$, proved from one axiom
  (the explicit formula).
\item The \textbf{det is $C^\infty$} is proved with linftyOp norm
  via the Leibniz formula (was axiom, now theorem).
\item The thimble geometry, Hamilton-Jacobi theory, and solution
  theory are documented but axiomatized. The
  \texttt{quantum\_thimble\_exists} axiom packages the IFT
  existence + BL variance bound.
\item The diamagnetic inequality has a proof structure
  (\texttt{DiagmagneticInequality.lean}, 413 lines) with the
  phase bound and diagonal case proved; the sub-steps (Trotter,
  heat kernel, Laplace transform) remain axioms (see
  \S\ref{sec:diamagnetic}).
\end{itemize}

\begin{thebibliography}{99}
\bibitem{BL76} H.~J.~Brascamp and E.~H.~Lieb,
  J.~Funct.~Anal.~\textbf{22} (1976) 366--389.
\bibitem{CD24} J.~G.~Conlon and J.~Dabkowski,
  ``Extensions of the Brascamp-Lieb inequality and the dipole gas,''
  arXiv:2412.05122 (2024).
\bibitem{DG25} P.~Dario and C.~Garban,
  Comm.~Math.~Phys.~(2025); arXiv:2311.16546.
\bibitem{Fukuma20} M.~Fukuma and N.~Matsumoto,
  ``Worldvolume approach to the tempered Lefschetz thimble method,''
  PTEP (2021); arXiv:2012.08468.
\bibitem{Fukuma23} M.~Fukuma,
  ``Simplified algorithm for the Worldvolume HMC,''
  arXiv:2311.10663 (2023).
\bibitem{Fukuma25b} M.~Fukuma and N.~Namekawa,
  ``Worldvolume HMC for the Hubbard model,''
  arXiv:2507.23748 (2025).
\bibitem{Fukuma26} M.~Fukuma,
  ``Worldvolume HMC for lattice gauge theories,''
  arXiv:2603.28629 (2026).
\bibitem{GJS74} J.~Glimm, A.~Jaffe, and T.~Spencer,
  Ann.~Math.~\textbf{100} (1974) 585--632.
\bibitem{Kupiainen80} A.~Kupiainen,
  Comm.~Math.~Phys.~\textbf{73} (1980) 273--294.
\bibitem{Kupiainen80b} A.~Kupiainen,
  Comm.~Math.~Phys.~\textbf{74} (1980) 199--222.
\bibitem{Lin24} Y.-H.~Lin,
  ``The Bayes principle and Segal axioms for P($\ph$)$_2$,''
  arXiv:2403.12804 (2024).
\bibitem{ABBPL18} A.~Alexandru, P.~F.~Bedaque, H.~Lamm, and S.~Lawrence,
  ``Finite-density Monte Carlo calculations on sign-optimized manifolds,''
  Phys.~Rev.~D~\textbf{97} (2018) 094510; arXiv:1804.00697.
\bibitem{LY24} S.~Lawrence and S.~Yamauchi,
  ``Convex optimization and contour deformations,''
  Phys.~Rev.~D~\textbf{110} (2024); arXiv:2311.13002.
\bibitem{CRS12} M.~Cristoforetti, F.~Di~Renzo, and L.~Scorzato,
  ``New approach to the sign problem in quantum field theories:
  High density QCD on a Lefschetz thimble,''
  Phys.~Rev.~D~\textbf{86} (2012) 074506; arXiv:1205.3996.
\bibitem{Scorzato15} L.~Scorzato,
  ``The Lefschetz thimble and the sign problem,''
  PoS LATTICE2015 (2016) 015; arXiv:1512.08039.
\bibitem{Simon1979} B.~Simon,
  \emph{Functional Integration and Quantum Physics},
  Academic Press (1979); 2nd ed., AMS Chelsea (2005).
\bibitem{TNK15} Y.~Tanizaki, H.~Nishimura, and K.~Kashiwa,
  ``Evading the sign problem in the mean-field approximation
  through Lefschetz-thimble path integral,''
  Phys.~Rev.~D~\textbf{91} (2015) 101701; arXiv:1504.02979.
\bibitem{Witten10} E.~Witten,
  ``A new look at the path integral of quantum mechanics,''
  in \emph{Surveys in Differential Geometry XV} (2010); arXiv:1009.6032.
\end{thebibliography}

\end{document}
